%% Generated by Sphinx.
\def\sphinxdocclass{report}
\documentclass[letterpaper,10pt,english]{sphinxmanual}
\ifdefined\pdfpxdimen
   \let\sphinxpxdimen\pdfpxdimen\else\newdimen\sphinxpxdimen
\fi \sphinxpxdimen=.75bp\relax
\ifdefined\pdfimageresolution
    \pdfimageresolution= \numexpr \dimexpr1in\relax/\sphinxpxdimen\relax
\fi
%% let collapsable pdf bookmarks panel have high depth per default
\PassOptionsToPackage{bookmarksdepth=5}{hyperref}

\PassOptionsToPackage{warn}{textcomp}


\usepackage{cmap}
\usepackage{fontspec}
\defaultfontfeatures[\rmfamily,\sffamily,\ttfamily]{}
\usepackage{amsmath,amssymb,amstext}
\usepackage{polyglossia}
\setmainlanguage{english}



        \setmainfont{Symbola}
    


\usepackage[Bjarne]{fncychap}
\usepackage{sphinx}

\fvset{fontsize=\small}
\usepackage{geometry}


% Include hyperref last.
\usepackage{hyperref}
% Fix anchor placement for figures with captions.
\usepackage{hypcap}% it must be loaded after hyperref.
% Set up styles of URL: it should be placed after hyperref.
\urlstyle{same}

\addto\captionsenglish{\renewcommand{\contentsname}{Getting Started}}

\usepackage{sphinxmessages}
\setcounter{tocdepth}{1}


        \usepackage{mathtools}
        \usepackage{newunicodechar}
        \usepackage[UTF8]{ctex}
        \newunicodechar{π}{\ensuremath{\pi}}
        \newunicodechar{≤}{\ensuremath{\le}}
        \newunicodechar{≥}{\ensuremath{\ge}}
        \newunicodechar{♥}{\ensuremath{\heartsuit}}
        \newunicodechar{…}{\ensuremath{\ldots}}
    

\title{BioModME}
\date{Jun 19, 2023}
\release{1.0.0}
\author{Justin Womack}
\newcommand{\sphinxlogo}{\vbox{}}
\renewcommand{\releasename}{Release}
\makeindex
\begin{document}

\pagestyle{empty}
\sphinxmaketitle
\pagestyle{plain}
\sphinxtableofcontents
\pagestyle{normal}
\phantomsection\label{\detokenize{index::doc}}



\chapter{Accessing BioModME}
\label{\detokenize{pages/INTRO_INDEX:accessing-biomodme}}\label{\detokenize{pages/INTRO_INDEX::doc}}
\sphinxAtStartPar
BioModME was built using the Rshiny framework with the injection of customized
HTML/CSS and javascript. As such, it is availble as a fully deployed web
application hosted by CTSI of the Medical College of Wisconsin. Additionally,
source code is provided for the application to be run locally.

\begin{sphinxadmonition}{note}{Note:}
\sphinxAtStartPar
Online Server is a good option that requires no download of application or
knowledge of running R.

\sphinxAtStartPar
We intend to package the application with Electron to have a fully
deployable version.
\end{sphinxadmonition}


\section{Use Online Server}
\label{\detokenize{pages/gettingStarted/useonline:use-online-server}}\label{\detokenize{pages/gettingStarted/useonline::doc}}
\sphinxAtStartPar
BioModME is hosted on webservers at the Medical College of Wiconsin by the
Clinical \& Translational Science Institute (CTSI) of Southeast Wisconsin.
It can be accessed at:

\sphinxAtStartPar
\sphinxurl{https://biomodme.ctsi.mcw.edu/}


\section{Download Locally}
\label{\detokenize{pages/gettingStarted/downloadtoR:download-locally}}\label{\detokenize{pages/gettingStarted/downloadtoR::doc}}
\sphinxAtStartPar
BioModME can be run on your local system using R.


\subsection{Install R and RStudio}
\label{\detokenize{pages/gettingStarted/downloadtoR:install-r-and-rstudio}}
\sphinxAtStartPar
R and RStudio will need to be installed (tested locally with R 4.1.1 and
RStudio 1.4.1717). Please check CRAN for the installation in R (
\sphinxurl{https://cran.r-project.org}).

\sphinxAtStartPar
For the installation of RStudio, go to and follow
instructions from \sphinxurl{https://www.rstudio.com}.


\subsection{Download BioModME from Github}
\label{\detokenize{pages/gettingStarted/downloadtoR:download-biomodme-from-github}}
\sphinxAtStartPar
\sphinxstylestrong{Option 1} \sphinxhyphen{} Using Git Bash:
\begin{enumerate}
\sphinxsetlistlabels{\arabic}{enumi}{enumii}{}{.}%
\item {} 
\sphinxAtStartPar
Open Git Bash

\item {} 
\sphinxAtStartPar
Change the current working directory to the location where you want the
cloned directory.

\item {} 
\sphinxAtStartPar
Enter the following command
\begin{quote}

\begin{sphinxVerbatim}[commandchars=\\\{\}]
\PYG{n}{git} \PYG{n}{clone} \PYG{n}{https}\PYG{p}{:}\PYG{o}{/}\PYG{o}{/}\PYG{n}{github}\PYG{o}{.}\PYG{n}{com}\PYG{o}{/}\PYG{n}{MCWComputationalBiologyLab}\PYG{o}{/}\PYG{n}{BioModME}\PYG{o}{.}\PYG{n}{git}
\end{sphinxVerbatim}
\end{quote}

\item {} 
\sphinxAtStartPar
Press \sphinxstylestrong{Enter} to create your local clone of the application.

\end{enumerate}

\sphinxAtStartPar
\sphinxstylestrong{Option 2} \sphinxhyphen{} Direct Download:
\begin{enumerate}
\sphinxsetlistlabels{\arabic}{enumi}{enumii}{}{.}%
\item {} 
\sphinxAtStartPar
Navigate to \sphinxurl{https://github.com/MCWComputationalBiologyLab/BioModME}.

\item {} 
\sphinxAtStartPar
Click the green button “<> Code”.

\noindent\sphinxincludegraphics{{download_code_button}.png}

\item {} 
\sphinxAtStartPar
Press download ZIP

\noindent{\hspace*{\fill}\sphinxincludegraphics[width=336\sphinxpxdimen,height=347\sphinxpxdimen]{{zip_download}.png}\hspace*{\fill}}

\item {} 
\sphinxAtStartPar
Unzip folder in desired location.

\end{enumerate}


\subsection{Install R Repositories}
\label{\detokenize{pages/gettingStarted/downloadtoR:install-r-repositories}}
\sphinxAtStartPar
BioModME uses a vast number of standard R repositories that will need to be
downloaded. Open up the UI.R file from the downloaded file. This script
loads all the needed libraries for the application to properly run.  To install
these repositories you can run the following:

\begin{sphinxVerbatim}[commandchars=\\\{\}]
\PYGZsh{} Vector of package names to install
load.lib \PYGZlt{}\PYGZhy{} c(\PYGZdq{}shinydashboard\PYGZdq{}, \PYGZdq{}bs4Dash\PYGZdq{}, \PYGZdq{}shiny\PYGZdq{}, \PYGZdq{}ggplot2\PYGZdq{}, \PYGZdq{}gridExtra\PYGZdq{},\PYGZdq{}shinythemes\PYGZdq{},
            \PYGZdq{}shinyWidgets\PYGZdq{}, \PYGZdq{}shinyjs\PYGZdq{}, \PYGZdq{}DT\PYGZdq{}, \PYGZdq{}tidyverse\PYGZdq{}, \PYGZdq{}dplyr\PYGZdq{}, \PYGZdq{}rhandsontable\PYGZdq{}, \PYGZdq{}data.table\PYGZdq{},
            \PYGZdq{}ggpmisc\PYGZdq{}, \PYGZdq{}colourpicker\PYGZdq{}, \PYGZdq{}shinyBS\PYGZdq{}, \PYGZdq{}shinyjqui\PYGZdq{}, \PYGZdq{}bsplus\PYGZdq{}, \PYGZdq{}deSolve\PYGZdq{}, \PYGZdq{}plotly\PYGZdq{},
            \PYGZdq{}Deriv\PYGZdq{}, \PYGZdq{}viridis\PYGZdq{}, \PYGZdq{}ggpubr\PYGZdq{}, \PYGZdq{}shinycssloaders\PYGZdq{}, \PYGZdq{}waiter\PYGZdq{}, \PYGZdq{}fresh\PYGZdq{}, \PYGZdq{}readxl\PYGZdq{},
            \PYGZdq{}minpack.lm\PYGZdq{}, \PYGZdq{}measurements\PYGZdq{}, \PYGZdq{}qdapRegex\PYGZdq{})

\PYGZsh{} Remove any packages from list that are already installed
install.lib \PYGZlt{}\PYGZhy{} load.lib[!load.lib \PYGZpc{}in\PYGZpc{} installed.packages()]

\PYGZsh{} Install all packages in install list
for (lib in install.lib) install.packages(lib, dependencies = TRUE)
sapply(load.lib, require, character = TRUE)
\end{sphinxVerbatim}

\sphinxAtStartPar
All the required packages should now be installed.


\subsection{Start Application}
\label{\detokenize{pages/gettingStarted/downloadtoR:start-application}}
\sphinxAtStartPar
The application can be started from base R by navigating to the application
home folder and running the following:

\begin{sphinxVerbatim}[commandchars=\\\{\}]
\PYG{n}{shiny}\PYG{p}{:}\PYG{p}{:}\PYG{n}{runApp}\PYG{p}{(}\PYG{p}{)}
\end{sphinxVerbatim}

\sphinxAtStartPar
A more common way to run the app is to run it directly from RStudio. This is
done by:
\begin{enumerate}
\sphinxsetlistlabels{\arabic}{enumi}{enumii}{}{.}%
\item {} 
\sphinxAtStartPar
Open either “ui.R” or “server.R” in RStudio

\item {} 
\sphinxAtStartPar
Press bottom in top right of script that says “Run App”.

\end{enumerate}

\begin{figure}[htbp]
\centering

\noindent\sphinxincludegraphics{{rstudio_runapp}.png}
\end{figure}


\section{Pull Docker Image}
\label{\detokenize{pages/gettingStarted/docker:pull-docker-image}}\label{\detokenize{pages/gettingStarted/docker::doc}}
\sphinxAtStartPar
The application has been created as a docker image for user convience if they
wish to run the app locally without installing R or worrying about other
possible dependencies.


\subsection{Download Docker}
\label{\detokenize{pages/gettingStarted/docker:download-docker}}
\sphinxAtStartPar
In this section, we will download docker to our computer.
This will install the desktop app but what we really need is the docker
terminal commands to be added.
If you have a working version of Docker on your system,
skip to the next section.
\begin{enumerate}
\sphinxsetlistlabels{\arabic}{enumi}{enumii}{}{.}%
\item {} 
\sphinxAtStartPar
Go to \sphinxurl{http://www.docker.com/products/docker-desktop}

\item {} 
\sphinxAtStartPar
Download docker for your respective system.
This tutorial follows the windows download.

\noindent\sphinxincludegraphics{{download_page}.png}

\item {} 
\sphinxAtStartPar
Go to you downloads folder and begin the system install of docker.
This download will take a few minutes.

\item {} 
\sphinxAtStartPar
A configuration popup will appear. Keep “Install required Windows components
for WSL 2” checked.  You can uncheck the desktop shortcut option if you want.

\noindent\sphinxincludegraphics{{download_options}.png}

\item {} 
\sphinxAtStartPar
Once download is complete a message will appear to restart your system.
Perform computer restart.

\item {} 
\sphinxAtStartPar
On restart you may get a notification that “WSL 2 installation is
incomplete”. If this is the case you will need to install this update or
the docker image will not run. Follow the link on the popup, it takes you
to the following link:
\sphinxurl{https://docs.microsoft.com/en-us/windows/wsl/install-manual\#step-4}–
\sphinxhyphen{}download\sphinxhyphen{}the\sphinxhyphen{}linux\sphinxhyphen{}kernel\sphinxhyphen{}update\sphinxhyphen{}package\}\{\sphinxurl{https://docs.microsoft.com/en-us}
/windows/wsl/install\sphinxhyphen{}manual\#step\sphinxhyphen{}4—download\sphinxhyphen{}the\sphinxhyphen{}linux\sphinxhyphen{}kernel\sphinxhyphen{}update\sphinxhyphen{}
package. Download and install the kernel from this link.

\noindent\sphinxincludegraphics{{download_wsl_popup}.png}

\item {} 
\sphinxAtStartPar
Restart computer again and open docker program.
It should be fully installed and working now.
A working program has a green color in the bottom left corner and lacks a
huge error message in the center of the application.

\noindent\sphinxincludegraphics{{working_docker}.png}

\end{enumerate}

\sphinxAtStartPar
Note: A Docker account does not need to be made in order to pull and run this
program.


\subsection{Pull/Run From Docker}
\label{\detokenize{pages/gettingStarted/docker:pull-run-from-docker}}
\sphinxAtStartPar
With docker installed, we can now use the command line to pull the program we
desire, meaning we can download it to our computer from the docker cloud.
\begin{enumerate}
\sphinxsetlistlabels{\arabic}{enumi}{enumii}{}{.}%
\item {} 
\sphinxAtStartPar
Open the computer’s terminal (command prompt).

\item {} 
\sphinxAtStartPar
Type the below command to pull the application from the docker cloud (this
takes approximately five minutes):

\begin{sphinxVerbatim}[commandchars=\\\{\}]
\PYG{n}{docker} \PYG{n}{pull} \PYG{n}{jwomack7512}\PYG{o}{/}\PYG{n}{biomodme}
\end{sphinxVerbatim}

\item {} 
\sphinxAtStartPar
Type the following command to run the application:
\begin{quote}

\begin{sphinxVerbatim}[commandchars=\\\{\}]
\PYG{n}{docker} \PYG{n}{run} \PYG{o}{\PYGZhy{}}\PYG{n}{dp} \PYG{l+m+mi}{3838}\PYG{p}{:}\PYG{l+m+mi}{3838} \PYG{n}{jwomack7512}\PYG{o}{/}\PYG{n}{biomodme}
\end{sphinxVerbatim}

\sphinxAtStartPar
This step might need approval from windows defender or other antivirus
program.

\noindent\sphinxincludegraphics{{terminal_commands}.png}
\end{quote}

\item {} 
\sphinxAtStartPar
The above command runs the program at a loack host port. To access
the app, open your web broswer and type the following:

\begin{sphinxVerbatim}[commandchars=\\\{\}]
\PYG{n}{localhost}\PYG{p}{:}\PYG{l+m+mi}{3838}
\end{sphinxVerbatim}

\item {} 
\sphinxAtStartPar
The application should now load on the opened webpage.

\end{enumerate}


\subsection{Loading Application After Pull}
\label{\detokenize{pages/gettingStarted/docker:loading-application-after-pull}}
\sphinxAtStartPar
After the image has been pulled once, there is not a need to repull it.
The image should be saved to your local docker program and can be loaded
easily using the Docker Hub GUI.
\begin{enumerate}
\sphinxsetlistlabels{\arabic}{enumi}{enumii}{}{.}%
\item {} 
\sphinxAtStartPar
Open docker and navigate to the “images” tab. Here you should see a copy
of the image you pulled in the previous step. This means its available
without having to repull it.

\item {} 
\sphinxAtStartPar
Scroll over the image \sphinxstylestrong{jwomack7512/biomodme}. A “run” button will appear.
Press that button.

\noindent\sphinxincludegraphics{{docker_hub_run_button}.png}

\item {} 
\sphinxAtStartPar
A popup will appear.
Click the dropdown arrow to show optional settings.
Here we give the container the name “bioModMe”.
If you do not give it a name, docker will randomly generate one.
Type in “3838” in Local Host.
This is necessary or the program will not load to the proper port the
Rshiny app specifies. Press the “Run” button.

\noindent{\hspace*{\fill}\sphinxincludegraphics[scale=0.7]{{load_container_options}.png}\hspace*{\fill}}

\item {} 
\sphinxAtStartPar
Go to the “Containers/Apps” tab in Docker Hub.
You should see the container you created with a note saying it is running
at port 3838.
You can access the app by simply clicking the “Open in Browser” button.
Note: you may already have a container built from using the command line,
meaning you could just run that one instead of rebuilding it in the previous
steps.

\noindent\sphinxincludegraphics{{docker_hub_container_to_terminal}.png}

\end{enumerate}

\sphinxAtStartPar
In the above figure, if you double click the container the logs will open for
the application. These can be useful for debugging and submitting errors.

\noindent\sphinxincludegraphics{{logs}.png}


\chapter{Page Layout}
\label{\detokenize{pages/API_INDEX:page-layout}}\label{\detokenize{pages/API_INDEX::doc}}

\section{Create Model}
\label{\detokenize{pages/api/CreateModel:create-model}}\label{\detokenize{pages/api/CreateModel::doc}}
\sphinxAtStartPar
General information for building a model.


\section{Execute Model}
\label{\detokenize{pages/api/ExecuteModel:execute-model}}\label{\detokenize{pages/api/ExecuteModel::doc}}
\sphinxAtStartPar
General information for executing the model.

\sphinxAtStartPar
new test for new stuff


\section{Visualization}
\label{\detokenize{pages/api/Visualization:visualization}}\label{\detokenize{pages/api/Visualization::doc}}
\sphinxAtStartPar
Holds plotting info, types, loading data, etc.


\chapter{Mathematical Derivations}
\label{\detokenize{pages/BACKEND_MATH_INDEX:mathematical-derivations}}\label{\detokenize{pages/BACKEND_MATH_INDEX::doc}}
\sphinxAtStartPar
This section covers all the mathematical equations and derivations that run in
the background of BioModME.


\section{Equations}
\label{\detokenize{pages/backend/eqn_info_index:equations}}\label{\detokenize{pages/backend/eqn_info_index::doc}}
\sphinxAtStartPar
This section contains the mathematical derivations used for the equation
builders found in this program.


\subsection{Chemical Based Equations}
\label{\detokenize{pages/backend/eqn_info/1_chem_eqns:chemical-based-equations}}\label{\detokenize{pages/backend/eqn_info/1_chem_eqns::doc}}

\subsubsection{Law of Mass Action}
\label{\detokenize{pages/backend/eqn_info/1_chem_eqns:law-of-mass-action}}
\sphinxAtStartPar
Chemical reactions are derived using the law of mass action. Given the following
chemical reaction scheme:
\begin{equation*}
\begin{split}\begin{equation*}
    aA + bB \longleftrightarrow[k_{-1}]{k_1} cC + dD
\end{equation*}\end{split}
\end{equation*}
\sphinxAtStartPar
where a, b, c, d are stoichiometic coefficients and A, B, C, D are the chemical
species. The law is as follows:
\begin{equation*}
\begin{split}\begin{equation*}
    -\frac{1}{a}\frac{d[A]}{dt} = -\frac{1}{b}\frac{d[B]}{dt} =
    \frac{1}{c}\frac{d[C]}{dt} = \frac{1}{d}\frac{d[D]}{dt} =
     k_1[A]^a[B]^b - k_{-1}[C]^c[D]^d
\end{equation*}\end{split}
\end{equation*}
\sphinxAtStartPar
From the above law, each individual term can is derived below:
\begin{equation*}
\begin{split}\begin{align*}
    \frac{d[A]}{dt} &= -a*k_1[A]^a[B]^b + a*k_2[C]^c[D]^d \\
    \frac{d[B]}{dt} &= -b*k_1[A]^a[B]^b + b*k_2[C]^c[D]^d \\
    \frac{d[C]}{dt} &= c*k_1[A]^a[B]^b - c*k_2[C]^c[D]^d \\
    \frac{d[D]}{dt} &= d*k_1[A]^a[B]^b - d*k_2[C]^c[D]^d
\end{align*}\end{split}
\end{equation*}
\sphinxAtStartPar
Each individual equation will be added to their overall flux differential
equation in the differential equation solver section.


\subsubsection{Law of Mass Action with Regulation}
\label{\detokenize{pages/backend/eqn_info/1_chem_eqns:law-of-mass-action-with-regulation}}
\sphinxAtStartPar
Mass action with regulation is for chemical reactions who have rate constants
that are dependent on other species. An example is the regulation of inactive
MPF to active MPF in the cell cycle by WEE1 and CDC25Cp. In this example,
inactive MPF is phosphorylated by CDC25Cp and dephosphorylated by WEE1. These
regulators are not directly affected in the their differential equations in
this situation while inactive and active MPF are. We have the same derivation
as mass action except the rate constants are changed to reflect the regulation.

\begin{figure}[htbp]
\centering

\noindent\sphinxincludegraphics[scale=0.4]{{mass_action_w_regulation}.png}
\end{figure}

\sphinxAtStartPar
The regulators will take the respective place of the rateconstants in the
mass action derivation.  In the mass action derivation above for “a” and “b”
we have:
\begin{equation*}
\begin{split}\begin{align*}
    \frac{d[A]}{dt} &= -a*k_1[A]^a[B]^b + a*k_2[C]^c[D]^d \\
    \frac{d[B]}{dt} &= -b*k_1[A]^a[B]^b + b*k_2[C]^c[D]^d
\end{align*}\end{split}
\end{equation*}
\sphinxAtStartPar
Adding regulator replaces the respective rate constant with the sum of
its regulators and rateconstants:
\begin{equation*}
\begin{split}\begin{align*}
    \frac{d[A]}{dt} &= -a*(\sum{[kf_i*reg_i]})[A]^a[B]^b +
    a*(\sum{[kr_i*reg_i]})[C]^c[D]^d \\
    \frac{d[B]}{dt} &= -b*(\sum{[kf_i*reg_i]})[A]^a[B]^b +
    b*(\sum{[kr_i*reg_i]})[C]^c[D]^d
\end{align*}\end{split}
\end{equation*}
\sphinxAtStartPar
Using the above example, we get the following result:
\begin{equation*}
\begin{split}\begin{align*}
    \frac{d[MPF]}{dt} &= -(k_{f1.1}*CDC25Cp)*MPF - (k_{r1.1}*WEE1)*MPFp \\
    \frac{d[MPFp]}{dt} &= (k_{f1.1}*CDC25Cp)*MPF + (k_{r1.1}*WEE1)*MPFp
\end{align*}\end{split}
\end{equation*}

\subsection{Enzyme Based Equations}
\label{\detokenize{pages/backend/eqn_info/2_enzyme_eqns:enzyme-based-equations}}\label{\detokenize{pages/backend/eqn_info/2_enzyme_eqns::doc}}

\subsubsection{Michaelis Menten}
\label{\detokenize{pages/backend/eqn_info/2_enzyme_eqns:michaelis-menten}}
\sphinxAtStartPar
Michaelis Menten kinetics is the most common approximation for describing the
rate at which enzyme reactions occur. The model most commonly takes the
following for describing the velocity of a reaction:
\begin{equation*}
\begin{split}\begin{equation*}
    v = \frac{d[P]}{dt} = V_{max}\frac{[S]}{K_M+[S]} =
    (k_{cat}*[E])\frac{[S]}{K_M+[S]}
\end{equation*}\end{split}
\end{equation*}
\sphinxAtStartPar
where,
\begin{quote}\begin{description}
\item[{\(V_{max}\)}] \leavevmode
\sphinxAtStartPar
Maximum Velocity

\item[{E}] \leavevmode
\sphinxAtStartPar
Enzyme Concentration

\item[{S}] \leavevmode
\sphinxAtStartPar
Substrate Concentration

\item[{P}] \leavevmode
\sphinxAtStartPar
Product Concentration

\item[{\(K_M\)}] \leavevmode
\sphinxAtStartPar
Michaelis Menten Constant

\item[{\(k_{cat}\)}] \leavevmode
\sphinxAtStartPar
Catalytic Rate of Enzyme Reaction

\end{description}\end{quote}


\subsection{Synthesis Reactions}
\label{\detokenize{pages/backend/eqn_info/3_synthesis:synthesis-reactions}}\label{\detokenize{pages/backend/eqn_info/3_synthesis::doc}}

\subsubsection{Synthesis}
\label{\detokenize{pages/backend/eqn_info/3_synthesis:synthesis}}
\sphinxAtStartPar
Synthesis by rate has a species generated at a specific rate. This is useful
for when the user does not particularly now what is causing the synthesis but
knows how quickly it is being produced. It simply creates a constant rate of
production.
\begin{equation*}
\begin{split}\begin{equation*}
    \frac{d[Species]}{dt} = k_{syn}
\end{equation*}\end{split}
\end{equation*}
\sphinxAtStartPar
where,
\begin{quote}\begin{description}
\item[{\(k_{syn}\)}] \leavevmode
\sphinxAtStartPar
synthesis rate of species

\end{description}\end{quote}


\subsubsection{Synthesis by Factor}
\label{\detokenize{pages/backend/eqn_info/3_synthesis:synthesis-by-factor}}
\sphinxAtStartPar
The “By Factor” law is used when a factor directly influences the synthesis of
a species but is not directly converted. An example of this is in the Cell
Cycle when the molecule E2F activates transcription factors that synthesize
Cyclin E and Cyclin A. If the cell contains no E2F, then there is no synthesis
of Cyclin E or A.

\begin{figure}[htbp]
\centering

\noindent\sphinxincludegraphics[width=400\sphinxpxdimen,height=200\sphinxpxdimen]{{synthesis_by_factor}.png}
\end{figure}

\sphinxAtStartPar
There is only one resultant mathematical flux from this type of synthesis and
that is the species being synthesized. The factor causing the synthesis does
not have its flux altered. The general resulting equations are:
\begin{equation*}
\begin{split}\begin{align*}
    \frac{d[Species]}{dt} &= k_{syn}*[Factor] \\
    \frac{[Factor]}{dt} &= 0
\end{align*}\end{split}
\end{equation*}
\sphinxAtStartPar
where,
\begin{quote}\begin{description}
\item[{\(k_{syn}\)}] \leavevmode
\sphinxAtStartPar
synthesis rate of species

\item[{Species}] \leavevmode
\sphinxAtStartPar
species that is being synthesized

\item[{Factor}] \leavevmode
\sphinxAtStartPar
factor that is causing the synthesis

\end{description}\end{quote}

\sphinxAtStartPar
In realtion to the above E2F equation, the differeneital equations are:
\begin{equation*}
\begin{split}\begin{align*}
    \frac{d[CycA]}{dt} &= k_{syn}*[E2F] \\
    \frac{d[CycE]}{dt} &= k_{syn}*[E2F] \\
    \frac{d[E2F]}{dt} &= 0 \\
\end{align*}\end{split}
\end{equation*}

\subsection{Degradation Reactions}
\label{\detokenize{pages/backend/eqn_info/4_degradation:degradation-reactions}}\label{\detokenize{pages/backend/eqn_info/4_degradation::doc}}

\subsubsection{Degradation}
\label{\detokenize{pages/backend/eqn_info/4_degradation:degradation}}
\sphinxAtStartPar
These reactions are dependant on a rate constant. Here we have two specific
options where the degradation can be concentration dependent or not.
Often degradations are concentration dependent on themselves, meaning there is
more degradation at higher concentrations and vice versa. The following is the
derivation for a concentration dependent degradation:
\begin{equation*}
\begin{split}\begin{equation*}
    \frac{d[Species]}{dt} = -k_{deg}*[Species]
\end{equation*}\end{split}
\end{equation*}
\sphinxAtStartPar
where,
\begin{quote}\begin{description}
\item[{\(k_{deg}\)}] \leavevmode
\sphinxAtStartPar
degradation rate constant

\item[{Species}] \leavevmode
\sphinxAtStartPar
species being degraded

\end{description}\end{quote}

\sphinxAtStartPar
When the degradation is not concentration dependent, the resulting flux is:
\begin{equation*}
\begin{split}\frac{d[Species]}{dt} = -k_{deg}\end{split}
\end{equation*}
\sphinxAtStartPar
When products are made they are just the opposite sign of the above reactions.
For example for concentration dependent reactions with one degradation to two
products we would have the resulting differential equations:
\begin{equation*}
\begin{split}\begin{align*}
    \frac{d[Species]}{dt} = -k_{deg}*[Species] \\
    \frac{d[Product 1]}{dt} = +k_{deg}*[Species] \\
    \frac{d[Product 2]}{dt} = +k_{deg}*[Species]
\end{align*}\end{split}
\end{equation*}

\subsubsection{Degradation by Enzyme}
\label{\detokenize{pages/backend/eqn_info/4_degradation:degradation-by-enzyme}}
\sphinxAtStartPar
This degradation option is a basic enzyme reaction using Michaelis Menten
kinetics where the substrate is the species to be degraded and their often
is no product. The resulting mathematical flux for the degraded substrate
would be:
\begin{equation*}
\begin{split}\frac{d[S]}{dt} = -V_{max}\frac{[S]}{K_m+[S]} =
-(k_{cat}*[E])\frac{[S]}{K_m+[S]}\end{split}
\end{equation*}
\sphinxAtStartPar
where,
\begin{quote}\begin{description}
\item[{\(V_{max}\)}] \leavevmode
\sphinxAtStartPar
Maximum Velocity

\item[{E}] \leavevmode
\sphinxAtStartPar
Enzyme Concentration

\item[{S}] \leavevmode
\sphinxAtStartPar
Substrate Concentration

\item[{P}] \leavevmode
\sphinxAtStartPar
Product Concentration

\item[{\(K_M\)}] \leavevmode
\sphinxAtStartPar
Michaelis Menten Constant

\item[{\(k_{cat}\)}] \leavevmode
\sphinxAtStartPar
Catalytic Rate of Enzyme Reaction

\end{description}\end{quote}

\sphinxAtStartPar
There is an option to add a product (P), which results in the mathematical
flux equations to break down to pure Michaelis Menten kinetics for the product
and substrate which would have the same equation above except with a positive
sign:
\begin{equation*}
\begin{split}\begin{align*}
    \frac{d[S]}{dt} &= -V_{max}\frac{[S]}{K_m+[S]} =
    -(k_{cat}*[E])\frac{[S]}{K_m+[S]} \\
    \frac{d[P]}{dt} &= +V_{max}\frac{[S]}{K_m+[S]} =
    +(k_{cat}*[E])\frac{[S]}{K_m+[S]}
\end{align*}\end{split}
\end{equation*}

\section{Input/Output}
\label{\detokenize{pages/backend/inputOutput_index:input-output}}\label{\detokenize{pages/backend/inputOutput_index::doc}}
\sphinxAtStartPar
This section contains the mathematical derivations used for the input/output
builders found in this program.


\subsection{Flow Between}
\label{\detokenize{pages/backend/io_info/1_flow_between:flow-between}}\label{\detokenize{pages/backend/io_info/1_flow_between::doc}}
\sphinxAtStartPar
Flow between two or more compartments follow the mechanics described below. In
the below figure, we have flow between two compartments. We have species
\(A_1\) flowing out of compartment 1 into compartment 2 as \(A_2\).

\begin{figure}[htbp]
\centering

\noindent\sphinxincludegraphics[width=400\sphinxpxdimen,height=300\sphinxpxdimen]{{flow_between}.png}
\end{figure}

\sphinxAtStartPar
The flows from a compartment are derived as the flow rate (F) multiplied by the
concentration of the species leaving the comparment as a minus term.
Flow to a compartment is the same derivation without the minus term.
The flows in the above diagram would be derived as:
\begin{equation*}
\begin{split}\begin{align*}
    V_1 * \frac{d[A_1]}{dt} &= -F * A_{1} \\
    V_2 * \frac{d[A_2]}{dt} &= F * A_{1}
\end{align*}\end{split}
\end{equation*}

\subsubsection{Split Flow}
\label{\detokenize{pages/backend/io_info/1_flow_between:split-flow}}
\sphinxAtStartPar
Often times we need to split the flow from one compartment to multiple
compartments.  This option splits the flow from one compartment to go to
multiple compartments. In the figure below, we see the flow from compartment 1
is split to go to compartment 2 and compartment 3.

\begin{figure}[htbp]
\centering

\noindent\sphinxincludegraphics[width=350\sphinxpxdimen,height=320\sphinxpxdimen]{{flow_between_split}.png}
\end{figure}

\sphinxAtStartPar
The summation of the output flows will be equal to the input flow:
\begin{equation*}
\begin{split}F_{out} = \sum_{1}^{n} F_{in}\end{split}
\end{equation*}
\sphinxAtStartPar
Given the above figure this equation derives out to:
\begin{equation*}
\begin{split}F = F_1 + F_2\end{split}
\end{equation*}
\sphinxAtStartPar
The resulting flow differential equations derive as:
\begin{equation*}
\begin{split}\begin{align*}
    V_{1} \frac{dA_{1}}{dt} = -F * A_{1} \\
    V_{2} \frac{dA_{2}}{dt} = -F_{1} * A_{1} \\
    V_{3} \frac{dA_{3}}{dt} = -F_{2} * A_{1} \\
\end{align*}\end{split}
\end{equation*}

\subsection{Flow In}
\label{\detokenize{pages/backend/io_info/2_flow_in:flow-in}}\label{\detokenize{pages/backend/io_info/2_flow_in::doc}}
\sphinxAtStartPar
Constant flow into a compartment from an unspecificed source. This could be
something like an input flow into a model of species or a known constant flow
from a source outside the scope of your created model. Flow In and Flow Out
options can also be pieced together to create a customized Flow Betweeen. In
the Figure below, we see a constant flow, \sphinxstylestrong{F}, carrying species, \sphinxstylestrong{A\_1},
into Compartment 1.

\begin{figure}[htbp]
\centering

\noindent\sphinxincludegraphics[width=350\sphinxpxdimen,height=100\sphinxpxdimen]{{flow_in}.png}
\end{figure}

\sphinxAtStartPar
The mathematical derivation for this model input would be:
\begin{equation*}
\begin{split}V_{1} \frac{dA_{1}}{dt} = F * A_{1}\end{split}
\end{equation*}

\subsection{Flow Out}
\label{\detokenize{pages/backend/io_info/3_flow_out:flow-out}}\label{\detokenize{pages/backend/io_info/3_flow_out::doc}}
\sphinxAtStartPar
Constant flow out of a compartment that does not go to another compartment.
Flow In and Flow Out options can be pieced together to create a customized Flow
Betweeen. In the Figure below, we see a constant flow, \sphinxstylestrong{F}, carrying species,
\sphinxstylestrong{A\_1}, out of Compartment 1.

\begin{figure}[htbp]
\centering

\noindent\sphinxincludegraphics[width=350\sphinxpxdimen,height=100\sphinxpxdimen]{{flow_out}.png}
\end{figure}

\sphinxAtStartPar
The mathematical derivation for this model output would be:
\begin{equation*}
\begin{split}V_{1} \frac{dA_{1}}{dt} = - F * A_{1}\end{split}
\end{equation*}

\subsection{Clearance}
\label{\detokenize{pages/backend/io_info/4_clearance:clearance}}\label{\detokenize{pages/backend/io_info/4_clearance::doc}}
\sphinxAtStartPar
Clearance from a compartment is used when a species leaves the compartment
and is removed from the model. In the example below, we have species
\(A_3\) being cleared from compartment 3.

\begin{figure}[htbp]
\centering

\noindent\sphinxincludegraphics[width=350\sphinxpxdimen,height=175\sphinxpxdimen]{{clearance}.png}
\end{figure}

\sphinxAtStartPar
It takes the following derivation:
\begin{equation*}
\begin{split}V_3 \frac{dA_3}{dt} = -k_e * A_3 * V_3\end{split}
\end{equation*}
\sphinxAtStartPar
where,
\begin{quote}\begin{description}
\item[{\(k_e\)}] \leavevmode
\sphinxAtStartPar
Rate of clearance

\item[{\(A_3\)}] \leavevmode
\sphinxAtStartPar
Concentration of Species being cleared

\item[{\(V_3\)}] \leavevmode
\sphinxAtStartPar
Volume of compartment that species is being removed from

\end{description}\end{quote}


\subsection{Simple Diffusion}
\label{\detokenize{pages/backend/io_info/5_simple_diffusion:simple-diffusion}}\label{\detokenize{pages/backend/io_info/5_simple_diffusion::doc}}
\sphinxAtStartPar
In simple diffusion, the substance is driven by concentration difference across
a membrane. In the example below, we have a liver compartment that is
separated into blood and liver compartments with a passive diffusion
membrance for compound \sphinxstylestrong{A}.

\begin{figure}[htbp]
\centering

\noindent\sphinxincludegraphics[width=400\sphinxpxdimen,height=200\sphinxpxdimen]{{simple_diffusion}.png}
\end{figure}

\sphinxAtStartPar
For the simple, one dimensional case shown above,
the mass flux through the membrane follows Ficks Law:
\begin{equation*}
\begin{split}J = -D \frac{dC}{dx}\end{split}
\end{equation*}
\sphinxAtStartPar
where,
\begin{quote}\begin{description}
\item[{C}] \leavevmode
\sphinxAtStartPar
Concentration at membrane location x

\item[{D}] \leavevmode
\sphinxAtStartPar
Diffusion coefficient of the molecule

\end{description}\end{quote}

\sphinxAtStartPar
We assume the concentration gradient across the membrane of thickness, x, falls
linearly with x leading to:
\begin{equation*}
\begin{split}\frac{dC}{dx} = \frac{\Delta C}{\Delta X} = \frac{C_2-C_1}{\Delta x}\end{split}
\end{equation*}\begin{equation*}
\begin{split}J = -D \frac{\Delta C}{\Delta x}\end{split}
\end{equation*}
\sphinxAtStartPar
The rate of transfer (\sphinxstylestrong{f}) of a neutral compound across a biological membrane of
surface area (\sphinxstylestrong{S}) is:
\begin{equation*}
\begin{split}f = JS = -SD\frac{\Delta C}{\Delta x} = -PS(C_2-C_1)\end{split}
\end{equation*}
\sphinxAtStartPar
where,
\begin{quote}\begin{description}
\item[{PS}] \leavevmode
\sphinxAtStartPar
Coefficient of diffusivity

\end{description}\end{quote}

\sphinxAtStartPar
Note: The \sphinxstylestrong{negative sign} on the right side of the equation indicates that
the net transfer due to diffusion is in a direction away from the region with
highter concentration.

\sphinxAtStartPar
Units on page:
\begin{description}
\item[{C}] \leavevmode
\sphinxAtStartPar
\(\text{conc}\)

\item[{D}] \leavevmode
\sphinxAtStartPar
\(\frac{\text{length}^2}{\text{time}}\)

\item[{PS}] \leavevmode
\sphinxAtStartPar
\(\frac{volume}{time}\)

\end{description}


\subsection{Facilitated Diffusion}
\label{\detokenize{pages/backend/io_info/6_facillitated_diffusion:facilitated-diffusion}}\label{\detokenize{pages/backend/io_info/6_facillitated_diffusion::doc}}
\sphinxAtStartPar
Facilitated  diffusion is a type of diffusion in which the molecules move from
the regions of higher concentration to the region of lower concentration
assisted by a carrier molecule. It is a selective process, only allowing
specific molecules to pass the gradient. In the below figure, we have a
liver compartment that is separated into blood and liver compartments with
a facilitated transportor for compound \sphinxstylestrong{A}. Note that this is a one way
transport.

\begin{figure}[htbp]
\centering

\noindent\sphinxincludegraphics[width=400\sphinxpxdimen,height=200\sphinxpxdimen]{{facillitated_diffusion}.png}
\end{figure}

\sphinxAtStartPar
For the derivation of the flux equation for facilitated diffusion, we take the
following equation:
\begin{equation*}
\begin{split}F  = \frac{V_{max}*[S_o]}{K_m}\end{split}
\end{equation*}
\sphinxAtStartPar
where,
\begin{quote}\begin{description}
\item[{\(V_{max}\)}] \leavevmode
\sphinxAtStartPar
Maximum Velocity

\item[{\(S_{o}\)}] \leavevmode
\sphinxAtStartPar
Substrate Concentration

\item[{\(K_{m}\)}] \leavevmode
\sphinxAtStartPar
Affinity of how much substrate concentration needed to reach \(V_{max}\)

\end{description}\end{quote}

\sphinxAtStartPar
The flux derivations for the above example would be:
\begin{equation*}
\begin{split}\begin{align*}
    V_{1} \frac{dA_1}{dt} &= - \frac{V_{max}*A_{1}}{k_m + A_1} \\
    V_{2} \frac{dA_2}{dt} &= \frac{V_{max}*A_{1}}{k_m + A_1} \\
\end{align*}\end{split}
\end{equation*}

\chapter{ODE Solvers}
\label{\detokenize{pages/backend/differential_solvers:ode-solvers}}\label{\detokenize{pages/backend/differential_solvers::doc}}
\sphinxAtStartPar
This section covers the differential solvers logic.


\chapter{Parameter Estimation}
\label{\detokenize{pages/backend/parameter_estimation:parameter-estimation}}\label{\detokenize{pages/backend/parameter_estimation::doc}}
\sphinxAtStartPar
This is the parameter estimation documentation.

\sphinxAtStartPar
Link to {\hyperref[\detokenize{pages/tutorial_par_estimation::doc}]{\sphinxcrossref{\DUrole{doc}{Parameter Estimation Tutorial}}}}


\chapter{Units}
\label{\detokenize{pages/backend/units:units}}\label{\detokenize{pages/backend/units::doc}}
\sphinxAtStartPar
This section will cover all thing units.  What units are allowed, how changing
units works, and how conversions are performed (maybe).


\chapter{Basic Tutorial}
\label{\detokenize{pages/tutorial_sc_index:basic-tutorial}}\label{\detokenize{pages/tutorial_sc_index::doc}}
\sphinxAtStartPar
This tutorial covers basics features of BioModME by building a simplified,
single compartment model. In this example, we focus on a set of reactions
occurring in a cell. Compound “A” is synthesized by “Prot” at a constant rate
and reacts with a limited supply of compound “B” found in the cell to make an
intermediate “C1”. C1 undergoes an enzymatic conversion to “C2” and this new
intermediate undergoes self\sphinxhyphen{}cleavage to the desired end product, “P”, and an
inhibitor, “I”, of Prot.  This inhibitor feedback binding to Prot stopping the
synthesis of A, and thereby the entire reaction scheme.

\begin{figure}[htbp]
\centering

\noindent\sphinxincludegraphics[width=0.500\linewidth]{{cell_model}.png}
\end{figure}

\begin{DUlineblock}{0em}
\item[] 
\end{DUlineblock}

\sphinxAtStartPar
In the following sections of this tutorial, we will design the above system and
look at how this program can be used to visualize and modify the data.
The contents include the following:
\begin{enumerate}
\sphinxsetlistlabels{\arabic}{enumi}{enumii}{}{.}%
\item {} 
\sphinxAtStartPar
Creating model variables

\item {} 
\sphinxAtStartPar
Building a system of equations

\item {} 
\sphinxAtStartPar
Entering the inputs and outputs of the system

\item {} 
\sphinxAtStartPar
Inputting numerical parameter values

\item {} 
\sphinxAtStartPar
Solving and plotting the differential model

\item {} 
\sphinxAtStartPar
Exploring the different types of plotting features

\item {} 
\sphinxAtStartPar
Exporting the model and its features in meaningful ways

\end{enumerate}

\sphinxAtStartPar
To use the online version of this application visit
\sphinxurl{https://biomodme.ctsi.mcw.edu/}.  Note this app is stored on an Rshiny server
that close unused applications after 15 minutes of inactivity. Your current
model can be downloaded in the “Export” tab as an .RDS file and can be loaded
in by opening the right sidebar of the application. The application can also
be directly downloaded at
\sphinxurl{https://github.com/MCWComputationalBiologyLab/BioModME} and opened using the
R programming language using stand shiny app protocol.


\section{Constructing Model}
\label{\detokenize{pages/tutorial_sc/1_construct_model:constructing-model}}\label{\detokenize{pages/tutorial_sc/1_construct_model::doc}}
\sphinxAtStartPar
This section will cover how to build and construct a model.


\subsection{Define Variables}
\label{\detokenize{pages/tutorial_sc/defineVariables:define-variables}}\label{\detokenize{pages/tutorial_sc/defineVariables::doc}}
\sphinxAtStartPar
The application loads on the “Create Model” tab.
Here, we will build our model, setting up variables, equations, and parameters.
Each of these components are their own box on the page.
On load, the application will create a starting compartment named “comp\_1”,
which can be seen in the “Compartments” box.
Click on the “Name” cell of the table to change the compartments name (cell).
We will also change the volume name here (V\_cell).
In a single compartment model, the physical volume is not of much importance
so we will leave that as a value of 1 to avoid any model scaling effects.

\begin{figure}[htbp]
\centering

\noindent\sphinxincludegraphics{{compartments}.png}
\end{figure}

\sphinxAtStartPar
Next, the model species will be added.  To add a species, move to the next box
(“Species”).  Here we press the button with a plus icon to add species whose
information can be edited in the table.

\begin{figure}[htbp]
\centering

\noindent\sphinxincludegraphics{{011_species_no_additions}.png}
\end{figure}

\sphinxAtStartPar
We will add nine species to the model and then rename them appropriately:
Prot, A, B, C\_1, C\_2, Enz, P, I, I.Prot

\begin{figure}[htbp]
\centering

\noindent\sphinxincludegraphics{{013_species_with_names}.png}
\end{figure}

\sphinxAtStartPar
The second column of the species is “Value”, which is the initial amount of
that species in the model.  For this model we have the following initial
concentrations:
\begin{itemize}
\item {} 
\sphinxAtStartPar
Prot = 3

\item {} 
\sphinxAtStartPar
A = 0

\item {} 
\sphinxAtStartPar
B = 5

\item {} 
\sphinxAtStartPar
C\_1 = 0

\item {} 
\sphinxAtStartPar
C\_2 = 0

\item {} 
\sphinxAtStartPar
Enz = 0.5

\item {} 
\sphinxAtStartPar
P = 0

\item {} 
\sphinxAtStartPar
I = 0

\item {} 
\sphinxAtStartPar
I.Prot = 0

\end{itemize}

\begin{figure}[htbp]
\centering

\noindent\sphinxincludegraphics{{014_species_IC}.png}
\end{figure}

\sphinxAtStartPar
Lastly, if desired, we can add descriptions to our species to remind ourselves
(or let other users) know what the species in the model are.  Below, we have
entered brief descriptions the beginning species. Editing can be done by
clicking the appropriate cell in the species table.

\begin{figure}[htbp]
\centering

\noindent\sphinxincludegraphics{{015_species_descriptions}.png}
\end{figure}


\subsection{Create Reactions}
\label{\detokenize{pages/tutorial_sc/createEquations:create-reactions}}\label{\detokenize{pages/tutorial_sc/createEquations::doc}}
\sphinxAtStartPar
Scroll down to the “Reactions” box in the “Create Model” tab. In this
section we will build up the sets of reactions that make up our reaction
network.  Press the “+” button to begin adding equations.


\subsubsection{Equation 1 \sphinxhyphen{} Law of Mass Action, Reversible}
\label{\detokenize{pages/tutorial_sc/createEquations:equation-1-law-of-mass-action-reversible}}
\noindent{\hspace*{\fill}\sphinxincludegraphics[scale=0.25]{{eqn1}.png}\hspace*{\fill}}

\begin{DUlineblock}{0em}
\item[] 
\end{DUlineblock}

\sphinxAtStartPar
Chemical reactions like this follow the \sphinxstylestrong{law of mass action}.
This first equation is a reversible reaction. The equation consists of two
reactants, one product, and moves in both directions where are forward rate
constant is “k\_f1” and our reverse is “k\_r1”.

\sphinxAtStartPar
This reaction input can be seen in the below figure.
The steps for the following equation are:
\begin{enumerate}
\sphinxsetlistlabels{\arabic}{enumi}{enumii}{}{.}%
\item {} 
\sphinxAtStartPar
Select “Chemical Reaction” for equation type.

\item {} 
\sphinxAtStartPar
Select “Mass Action” as the law.

\item {} 
\sphinxAtStartPar
Select “Reversible” as the reaction direction.

\item {} 
\sphinxAtStartPar
Set the equation to have two reactants and one product.

\item {} 
\sphinxAtStartPar
Fill out equation information. Rate constant names are autogenerated for all
equation types using the number of the current equation as their number
at the end.

\item {} 
\sphinxAtStartPar
This is the equation builder that shows what your built equation looks like.

\item {} 
\sphinxAtStartPar
Press the “Add Equation” button.

\end{enumerate}

\noindent{\hspace*{\fill}\sphinxincludegraphics{{02_reactions_eqn_1_marked}.png}\hspace*{\fill}}

\begin{DUlineblock}{0em}
\item[] 
\end{DUlineblock}


\subsubsection{Equation 2 \sphinxhyphen{} Michaelis\sphinxhyphen{}Menton Kinetics}
\label{\detokenize{pages/tutorial_sc/createEquations:equation-2-michaelis-menton-kinetics}}
\noindent{\hspace*{\fill}\sphinxincludegraphics[scale=0.15]{{eqn2}.png}\hspace*{\fill}}

\begin{DUlineblock}{0em}
\item[] 
\end{DUlineblock}

\sphinxAtStartPar
Next, we enter the second equation, conversion of C1 to C2 by the enzyme “Enz”.
These reactions follow Michaelis\sphinxhyphen{}Menten kinetics. Figure 7 shows how this
equation looks entered the program.

\sphinxAtStartPar
The steps are as follows:
\begin{enumerate}
\sphinxsetlistlabels{\arabic}{enumi}{enumii}{}{.}%
\item {} 
\sphinxAtStartPar
Change the Reaction Type dropdown to “Enzyme Based Reaction”.

\item {} 
\sphinxAtStartPar
Select “Michaelis Menten” as the Law.

\item {} 
\sphinxAtStartPar
Enter the equation information into the builder.

\item {} 
\sphinxAtStartPar
Press the “Add Equation” Button.

\end{enumerate}

\noindent{\hspace*{\fill}\sphinxincludegraphics{{02_reactions_eqn_2_marked}.png}\hspace*{\fill}}

\begin{DUlineblock}{0em}
\item[] 
\end{DUlineblock}


\subsubsection{Equation 3 \sphinxhyphen{} Law of Mass Action, Unidirection}
\label{\detokenize{pages/tutorial_sc/createEquations:equation-3-law-of-mass-action-unidirection}}
\noindent{\hspace*{\fill}\sphinxincludegraphics[scale=0.62]{{eqn3}.png}\hspace*{\fill}}

\begin{DUlineblock}{0em}
\item[] 
\end{DUlineblock}

\sphinxAtStartPar
The next reaction is a conversion of C2 to product P and inhibitor I.
This is a one directional reaction with one reactant and two products following
the law of mass action.
\begin{enumerate}
\sphinxsetlistlabels{\arabic}{enumi}{enumii}{}{.}%
\item {} 
\sphinxAtStartPar
Change the Equation Type dropdown to “Chemical Reaction”.

\item {} 
\sphinxAtStartPar
Select “Mass Action” as the Law.

\item {} 
\sphinxAtStartPar
Change the reversability option to “Irreversible”.

\item {} 
\sphinxAtStartPar
Set the Number of Reactants to “1” and the Number of Products to “2”.

\item {} 
\sphinxAtStartPar
Enter the equation information into the builder.

\item {} 
\sphinxAtStartPar
Press “Add Equation” button.

\end{enumerate}

\noindent{\hspace*{\fill}\sphinxincludegraphics{{02_reactions_eqn_3_marked}.png}\hspace*{\fill}}

\begin{DUlineblock}{0em}
\item[] 
\end{DUlineblock}


\subsubsection{Equation 4 \sphinxhyphen{} Law of Mass Action, Unidirection}
\label{\detokenize{pages/tutorial_sc/createEquations:equation-4-law-of-mass-action-unidirection}}
\noindent{\hspace*{\fill}\sphinxincludegraphics[scale=0.5]{{eqn4}.png}\hspace*{\fill}}

\begin{DUlineblock}{0em}
\item[] 
\end{DUlineblock}

\sphinxAtStartPar
Our next reaction is the inhibitor of Prot. This is another chemical reaction
following the law of mass as in Equations 1 and 3 with two Reactants I and Prot
converted to I.Prot. The user interface input is below in Figure 9.
The steps should be familiar to your by now, but they are as follows:
\begin{enumerate}
\sphinxsetlistlabels{\arabic}{enumi}{enumii}{}{.}%
\item {} 
\sphinxAtStartPar
Change the Reaction Type dropdown to “Chemical Reaction”.

\item {} 
\sphinxAtStartPar
Select “Mass Action” as the Law.

\item {} 
\sphinxAtStartPar
Change the reversability option to “Irreversible”.

\item {} 
\sphinxAtStartPar
Set the Number of Reactants to “2” and the Number of Products to “1”.

\item {} 
\sphinxAtStartPar
Enter the equation information into the builder.

\item {} 
\sphinxAtStartPar
Press “Add Equation” button.

\end{enumerate}

\noindent{\hspace*{\fill}\sphinxincludegraphics{{02_reactions_eqn_4_marked}.png}\hspace*{\fill}}

\begin{DUlineblock}{0em}
\item[] 
\end{DUlineblock}


\subsubsection{Equation 5 \sphinxhyphen{} Synthesis by Factor}
\label{\detokenize{pages/tutorial_sc/createEquations:equation-5-synthesis-by-factor}}
\noindent{\hspace*{\fill}\sphinxincludegraphics[scale=0.5]{{eqn5}.png}\hspace*{\fill}}

\begin{DUlineblock}{0em}
\item[] 
\end{DUlineblock}

\sphinxAtStartPar
Our next reaction is the synthesis of A by Prot, where Prot is a factor
promoting synthesis (as opposed to being directly converted).
\begin{enumerate}
\sphinxsetlistlabels{\arabic}{enumi}{enumii}{}{.}%
\item {} 
\sphinxAtStartPar
Set the Reaction Type to “Chemical Reaction”.

\item {} 
\sphinxAtStartPar
Change the Reaction Law dropdown to “Synthesis”.

\item {} 
\sphinxAtStartPar
Click the “Factor Driving Synthesis?” checbox.

\item {} 
\sphinxAtStartPar
Fill out equation builder with the species to synthesize and its
corresponding factor.

\item {} 
\sphinxAtStartPar
Press the “Add Equation” button

\end{enumerate}

\noindent{\hspace*{\fill}\sphinxincludegraphics{{02_reactions_eqn_5_marked}.png}\hspace*{\fill}}

\begin{DUlineblock}{0em}
\item[] 
\end{DUlineblock}


\subsubsection{Equation 6 \sphinxhyphen{} Degradation by Rate}
\label{\detokenize{pages/tutorial_sc/createEquations:equation-6-degradation-by-rate}}
\noindent{\hspace*{\fill}\sphinxincludegraphics[scale=0.5]{{eqn6}.png}\hspace*{\fill}}

\begin{DUlineblock}{0em}
\item[] 
\end{DUlineblock}

\sphinxAtStartPar
Our final reaction is the degradation of I.Prot by a rate.
This is useful for when we know the rate at which a protein is degraded in
the cell but do not really know what is causing the degradation.
These are often concentration dependent.
\begin{enumerate}
\sphinxsetlistlabels{\arabic}{enumi}{enumii}{}{.}%
\item {} 
\sphinxAtStartPar
Set the Reaction Type to “Chemical Reaction”.

\item {} 
\sphinxAtStartPar
Change the Reaction Law dropdown to “Degradation (Rate)”.

\item {} 
\sphinxAtStartPar
Fill out equation builder with the species to degrade and its rate constant.

\item {} 
\sphinxAtStartPar
Make sure the “Concentration Dependent” box is checked.

\item {} 
\sphinxAtStartPar
Press the “Add Equation” button.

\end{enumerate}

\noindent{\hspace*{\fill}\sphinxincludegraphics{{02_reactions_eqn_6_marked}.png}\hspace*{\fill}}

\begin{DUlineblock}{0em}
\item[] 
\end{DUlineblock}


\subsection{Compartment Input/Output}
\label{\detokenize{pages/tutorial_sc/compartmentIO:compartment-input-output}}\label{\detokenize{pages/tutorial_sc/compartmentIO::doc}}
\sphinxAtStartPar
The next box is the Input/Output box.  This box is used to add flow/transfers
between compartments.  Our model does not have any flows to account for so we
will not use this box for anything.

\noindent{\hspace*{\fill}\sphinxincludegraphics{{03_IO_box}.png}\hspace*{\fill}}

\begin{DUlineblock}{0em}
\item[] 
\end{DUlineblock}


\subsection{Parameters}
\label{\detokenize{pages/tutorial_sc/parameters:parameters}}\label{\detokenize{pages/tutorial_sc/parameters::doc}}
\sphinxAtStartPar
Navigate to the next box: \sphinxstylestrong{Parameters}.

\sphinxAtStartPar
Below is the generated parameter table that has been created so far by the
above entered information. All descriptions are pregenerated based on where
the parameter originated. Like other tables, all values are editable.

\begin{figure}[htbp]
\centering

\noindent\sphinxincludegraphics{{04_parameters}.png}
\end{figure}

\begin{DUlineblock}{0em}
\item[] 
\end{DUlineblock}

\sphinxAtStartPar
The parameter values for the model are as follows:
\begin{itemize}
\item {} 
\sphinxAtStartPar
k\_f1 = 2

\item {} 
\sphinxAtStartPar
k\_r1 = 1

\item {} 
\sphinxAtStartPar
Km\_2 = 0.1

\item {} 
\sphinxAtStartPar
kcat\_2 = 3

\item {} 
\sphinxAtStartPar
k\_f3 = 2

\item {} 
\sphinxAtStartPar
k\_f4 = 2

\item {} 
\sphinxAtStartPar
k\_s5 = 0.5

\item {} 
\sphinxAtStartPar
k\_d6 = 0.1

\end{itemize}

\sphinxAtStartPar
The below figure shows the parameter table with the values entered.

\begin{figure}[htbp]
\centering

\noindent\sphinxincludegraphics{{041_parameters_w_values}.png}
\end{figure}

\begin{DUlineblock}{0em}
\item[] 
\end{DUlineblock}


\subsection{Differential Equations}
\label{\detokenize{pages/tutorial_sc/diffeqns:differential-equations}}\label{\detokenize{pages/tutorial_sc/diffeqns::doc}}
\sphinxAtStartPar
Navigate to the final box “Differential Equations”.

\sphinxAtStartPar
This tab is used to generate and see the mathematical equations that govern
our model.  This tab will begin with all species being equal to “NA” until the
generate button is pressed or model is executed.

\sphinxAtStartPar
Differential equations are solved for whenever an equation is added.
The “Generate” button is used to refresh the viewing of the differential
equations if necessary.

\begin{figure}[htbp]
\centering

\noindent\sphinxincludegraphics{{05_differential_equations}.png}
\end{figure}

\sphinxAtStartPar
\sphinxstylestrong{Note:} The right sidebar of the “System of Differential Equations” box has an
option for viewing the differential equations including turning off the math
rendering and showing each derivation line as a newline for each differential
equation.


\subsection{Solve Model}
\label{\detokenize{pages/tutorial_sc/executemodel:solve-model}}\label{\detokenize{pages/tutorial_sc/executemodel::doc}}
\sphinxAtStartPar
Navigate to the “Execute Model” tab on the lefthand tabbar. This tab solves our
model and provides us with the options to customize the solver. For this model,
we will just solve using the standard options:
\begin{enumerate}
\sphinxsetlistlabels{\arabic}{enumi}{enumii}{}{.}%
\item {} 
\sphinxAtStartPar
Enter the following times: Starting = 0, End = 15, Step = 0.1, Unit = min.
This generates the model to be solved and integrated at time points 0, 0.1,
0.2…..14.8, 14.9, 15.

\item {} 
\sphinxAtStartPar
Press the “Run Solver” button.

\item {} 
\sphinxAtStartPar
An output table should be generated with the concentration of each species
along each time step of the solved model.

\end{enumerate}

\sphinxAtStartPar
You can download this solved model data to a csv file in the
download tab on this page.  Post processing is not used in this tutorial.

\begin{figure}[htbp]
\centering

\noindent\sphinxincludegraphics{{10_execute_marked}.png}
\end{figure}


\section{Visualize Model}
\label{\detokenize{pages/tutorial_sc/2_visualize_model:visualize-model}}\label{\detokenize{pages/tutorial_sc/2_visualize_model::doc}}
\sphinxAtStartPar
This section will cover how to visualize our created model.

\sphinxAtStartPar
Navigate to the “Visualization” tab on the lefthand tabbar.

\sphinxAtStartPar
The plotting tab has four modes, we will cover two of them. In this section we
will cover “Normal Plot” which is the default when you enter the tab.
The tab auto loads the plot when the model is solved, plotting all species
of the model. These are the main features of the plot:
\begin{enumerate}
\sphinxsetlistlabels{\arabic}{enumi}{enumii}{}{.}%
\item {} 
\sphinxAtStartPar
Variables \sphinxhyphen{} Changes what species are plotting in the model.
All species in model are plotted by default.

\item {} 
\sphinxAtStartPar
Download \sphinxhyphen{} Dropdown containing options to download plot.

\item {} 
\sphinxAtStartPar
Options \sphinxhyphen{} Dropdown containing the following:
\begin{itemize}
\item {} 
\sphinxAtStartPar
Plot mode \sphinxhyphen{} Changes the plotting types and features.
Here we are plotting “Normal Plot”.

\item {} 
\sphinxAtStartPar
Plot Renderer – type of plot used. “plotly” is an interactive plot while
“ggplot2” is the standard R plotting. The interactive is the default as it
allows you to interact with the data using features such as zoom, pan,
and hovering over data to view its value.

\end{itemize}

\item {} 
\sphinxAtStartPar
Plot \sphinxhyphen{} Plot area with adjustable sizing on bottom right corner.

\item {} 
\sphinxAtStartPar
Plot Options – Four dropdowns that allow the addition of labels,
change of axis sizing, and control of line color and size.

\item {} 
\sphinxAtStartPar
Boxes containing menus to change model variables and import datasets to
plot.

\end{enumerate}

\begin{figure}[htbp]
\centering

\noindent\sphinxincludegraphics{{20_plot_marked}.png}
\end{figure}

\sphinxAtStartPar
By default the page should look like above when this tutorial model is solved.
All variables are displayed using a evenly spaced rainbow pattern coloring.

\sphinxAtStartPar
Press the “Variables” dropdown to change the species plotted in the graph.
Press the “x” on the species to remove them.  The “Select All” button
reselects all species while the “Reset” button removes all species.

\begin{figure}[htbp]
\centering

\noindent\sphinxincludegraphics{{21_plot_var_button}.png}
\end{figure}

\sphinxAtStartPar
Let’s remove the following species to clean the plot up:
\begin{itemize}
\item {} 
\sphinxAtStartPar
C\_1

\item {} 
\sphinxAtStartPar
C\_2

\item {} 
\sphinxAtStartPar
Enz

\end{itemize}

\sphinxAtStartPar
We now have a less cluttered plot:

\begin{figure}[htbp]
\centering

\noindent\sphinxincludegraphics{{22_plot_less_cluttered}.png}
\end{figure}

\sphinxAtStartPar
The interactive plot uses Plotly (\sphinxurl{https://plotly.com/r/plotly-fundamentals/})
as a render and includes many of its base features including:
\begin{enumerate}
\sphinxsetlistlabels{\arabic}{enumi}{enumii}{}{.}%
\item {} 
\sphinxAtStartPar
By clicking on a variable in the legend, they can be removed from the plot.

\item {} 
\sphinxAtStartPar
Double clicking a legend variable will remove all species but that one
from the plot.

\item {} 
\sphinxAtStartPar
The top right bar includes options to pan, hover, and zoom the plot.

\item {} 
\sphinxAtStartPar
Press, hold, and drag the mouse on the plot to select specific features.

\item {} 
\sphinxAtStartPar
By hovering over lines on the plot, a popup will show with the values
at that point.

\end{enumerate}

\sphinxAtStartPar
There are many features included to let the user make the plot there own.
These are all located in the dropdowns below the plot.
Here, almost any feature of the plot is customizable.
Feel free to play around and make the plot your own.


\section{Export}
\label{\detokenize{pages/tutorial_sc/3_export:export}}\label{\detokenize{pages/tutorial_sc/3_export::doc}}
\sphinxAtStartPar
BioModME allows the built model to be exported into 4 main categories:
\begin{enumerate}
\sphinxsetlistlabels{\arabic}{enumi}{enumii}{}{.}%
\item {} 
\sphinxAtStartPar
Prebuild, ready to run code, for either MATLAB or R.

\item {} 
\sphinxAtStartPar
An “.rds” file (type of R data file) that can be reloaded into the program.
This is the way to save and load a model into the program currently.

\item {} 
\sphinxAtStartPar
A prebuild latex document all main model information.

\item {} 
\sphinxAtStartPar
Downloading tables for parameters, equations, and initial conditions in csv,
txt, or pdf form.

\end{enumerate}

\begin{figure}[htbp]
\centering

\noindent\sphinxincludegraphics{{30_export_page_marked}.png}
\end{figure}

\sphinxAtStartPar
Feel free to try the code export if you have an interest in solving your
model in those programming languages.  Make sure to give any downloaded file
a name.


\subsection{Latex}
\label{\detokenize{pages/tutorial_sc/3_export:latex}}
\sphinxAtStartPar
When downloading the latex document a popup will appear, showing
further customization options for that document.

\begin{figure}[htbp]
\centering

\noindent\sphinxincludegraphics{{31_export_latex_options}.png}
\end{figure}

\sphinxAtStartPar
In this example, we will remove the “Additional equations” page as we did not
create any non\sphinxhyphen{}standard equations. The generated document is created as a text
file that can be copied to our latex editor of choice and run.  I recommend
using overleaf (www.overleaf.com) but any TeX editor will work.
Below is a figure of the overleaf output:

\begin{figure}[htbp]
\centering

\noindent\sphinxincludegraphics{{overleaf}.png}
\end{figure}


\chapter{Multi\sphinxhyphen{}Compartmental Model}
\label{\detokenize{pages/tutorial_mc_index:multi-compartmental-model}}\label{\detokenize{pages/tutorial_mc_index::doc}}
\sphinxAtStartPar
A bolus of a neutral drug “A” was administrated intravenously (i.v.) at time
t = 0. Within the liver, “A” can diffuse (unidirectional) from liver blood
region tissue liver tissue region via facilitated diffusion with Vmax
(nmoles/min) and Km (Molar). Within the liver cells, “A” is converted to
“B” via a chemical reaction that follows the law of mass action with a first
order rate constant k1 (1/min). The product “B” can also diffuse
(simple diffusion) between liver tissue region and blood region with a
permeability\sphinxhyphen{}surface area product PS (ml/min). Within the kidneys, “B” is
excreted from the blood region via a linear process with a rate constant ke
(1/min).

\begin{figure}[htbp]
\centering

\noindent\sphinxincludegraphics[scale=0.4]{{flow_diagram}.png}
\end{figure}

\sphinxAtStartPar
The following are the parameter and concentration tables for the model:

\begin{figure}[htbp]
\centering

\noindent\sphinxincludegraphics{{information_tables}.png}
\end{figure}

\sphinxAtStartPar
In the following sections of this tutorial, we will design the above
multi compartment method.

\sphinxAtStartPar
The contents focus on:
\begin{enumerate}
\sphinxsetlistlabels{\arabic}{enumi}{enumii}{}{.}%
\item {} 
\sphinxAtStartPar
Adding multiple compartments.

\item {} 
\sphinxAtStartPar
Adding variables between compartments.

\item {} 
\sphinxAtStartPar
Properly assigning flows between compartments

\item {} 
\sphinxAtStartPar
Using transporters to pass species between compartments.

\end{enumerate}

\sphinxAtStartPar
To use the online version of this application visit
\sphinxurl{https://biomodme.ctsi.mcw.edu/}.  Note this app is stored on an Rshiny server
that close unused applications after 15 minutes of inactivity. Your current
model can be downloaded in the “Export” tab as an .RDS file and can be loaded
in by opening the right sidebar of the application. The application can also
be directly downloaded at
\sphinxurl{https://github.com/MCWComputationalBiologyLab/BioModME} and opened using the
R programming language using stand shiny app protocol.


\section{Constructing Model}
\label{\detokenize{pages/tutorial_mc/1_construct_model:constructing-model}}\label{\detokenize{pages/tutorial_mc/1_construct_model::doc}}
\sphinxAtStartPar
This section will cover how to build and construct a model.


\subsection{Create Compartments}
\label{\detokenize{pages/tutorial_mc/addCompartments:create-compartments}}\label{\detokenize{pages/tutorial_mc/addCompartments::doc}}
\sphinxAtStartPar
We will start by setting up our four compartments in this model:
Rest\_Of\_Body, Liver\_Blood, Liver\_Tissue, Kidney.

\sphinxAtStartPar
Steps:
\begin{enumerate}
\sphinxsetlistlabels{\arabic}{enumi}{enumii}{}{.}%
\item {} 
\sphinxAtStartPar
Go to the \sphinxstylestrong{Compartments} box in the \sphinxstylestrong{Create Model} tab.

\item {} 
\sphinxAtStartPar
Press the “\sphinxstylestrong{+}”” button three times to add three comparments.

\item {} 
\sphinxAtStartPar
Change the Names of the compartments in the table to \sphinxstylestrong{Rest\_Of\_Body},
\sphinxstylestrong{Liver\_Blood}, \sphinxstylestrong{Liver\_Tissue}, and \sphinxstylestrong{Kidney}.

\item {} 
\sphinxAtStartPar
Volume names are autogenerated with numerical numbering.

\item {} 
\sphinxAtStartPar
Change the volume value and units to match parameter table (these values can
also be changed in the parameter table for the model).

\end{enumerate}

\sphinxAtStartPar
The end result should be the following table:

\begin{figure}[htbp]
\centering

\noindent\sphinxincludegraphics{{comp_4_vol_values}.png}
\end{figure}


\subsection{Define Variables}
\label{\detokenize{pages/tutorial_mc/addSpecies:define-variables}}\label{\detokenize{pages/tutorial_mc/addSpecies::doc}}
\sphinxAtStartPar
There are two variables (\sphinxstylestrong{A} \& \sphinxstylestrong{B}) to add to all four comparments.
We will make use of the “Add To All Compartments” option to quickly add the
same species to all compartments.

\sphinxAtStartPar
Steps:
\begin{enumerate}
\sphinxsetlistlabels{\arabic}{enumi}{enumii}{}{.}%
\item {} 
\sphinxAtStartPar
Select the “Add To All Compartments” checkbox.

\begin{figure}[htbp]
\centering

\noindent\sphinxincludegraphics{{species_1_marked}.png}
\end{figure}

\item {} 
\sphinxAtStartPar
Press the “\sphinxstylestrong{+}” button to generate a popup.

\item {} 
\sphinxAtStartPar
In Popup type “A B” separated by a space. Check “Numerical” instead of
“Compartment Name” and press the add button. This will add the variables
\sphinxstylestrong{A\_1}, \sphinxstylestrong{A\_2}, \sphinxstylestrong{A\_3}, \sphinxstylestrong{A\_4}, \sphinxstylestrong{B\_1}, \sphinxstylestrong{B\_2}, \sphinxstylestrong{B\_3}, and \sphinxstylestrong{B\_4}
to the model.

\begin{figure}[htbp]
\centering

\noindent\sphinxincludegraphics{{species_2_enter}.png}
\end{figure}

\item {} 
\sphinxAtStartPar
Unselect the “Show Active Compartment Only Checkmark” to see all the species
in the model.

\begin{figure}[htbp]
\centering

\noindent\sphinxincludegraphics{{species_3_all}.png}
\end{figure}

\item {} 
\sphinxAtStartPar
Change the concentration of \sphinxstylestrong{A\_1} to match initial condition table.

\begin{figure}[htbp]
\centering

\noindent\sphinxincludegraphics{{species_4_conc_added}.png}
\end{figure}

\end{enumerate}


\subsection{Add Equations}
\label{\detokenize{pages/tutorial_mc/addEquations:add-equations}}\label{\detokenize{pages/tutorial_mc/addEquations::doc}}
\sphinxAtStartPar
This tutorial model has only one reaction: the conversion of A to B (\sphinxstylestrong{A3} \&
\sphinxstylestrong{B3}) in the liver tissue.

\begin{figure}[htbp]
\centering

\noindent\sphinxincludegraphics[scale=0.75]{{liver_equation}.png}
\end{figure}

\sphinxAtStartPar
Steps:
\begin{enumerate}
\sphinxsetlistlabels{\arabic}{enumi}{enumii}{}{.}%
\item {} 
\sphinxAtStartPar
Scroll down to the \sphinxstylestrong{Reactions} box in the \sphinxstylestrong{Create Model} tab. Press
the addition button below the equation table.

\item {} 
\sphinxAtStartPar
Change the \sphinxstylestrong{Active Compartment} to “Liver\_Tissue”.

\item {} 
\sphinxAtStartPar
Make sure \sphinxstylestrong{Mass Action} is selected as the \sphinxstylestrong{Law} for this reaction.
Change the \sphinxstylestrong{Reaction Direction} to “Forward” for this is a one way
reaction.

\item {} 
\sphinxAtStartPar
This reaction has one reactant and one product. Make sure \sphinxstylestrong{Reactant 1} is
\sphinxstylestrong{A\_3} and \sphinxstylestrong{Product 1} is \sphinxstylestrong{B\_3}. Make the \sphinxstylestrong{Forward Rate Constant}
to \sphinxstylestrong{k\_f1}, if not already. You can set the parameter value here or later
on the parameter table. For this tutorial most values will be set on the
parameter table.

\item {} 
\sphinxAtStartPar
Press the \sphinxstylestrong{Add Equation} button.

\end{enumerate}

\begin{figure}[htbp]
\centering

\noindent\sphinxincludegraphics{{eqn_1_marked}.png}
\end{figure}


\subsection{Add Compartment Inputs/Outputs}
\label{\detokenize{pages/tutorial_mc/addIO:add-compartment-inputs-outputs}}\label{\detokenize{pages/tutorial_mc/addIO::doc}}
\sphinxAtStartPar
In this section of the tutorial, we will cover how to add flows, inputs, and
outputs of compartments to the model. Begin by scroling to the
\sphinxstylestrong{Input/Output} box on the \sphinxstylestrong{Create Model} tab. Press the addition button to
begin adding the below flows between species information.
Again, here is the overall flow model:

\begin{figure}[htbp]
\centering

\noindent\sphinxincludegraphics[scale=0.4]{{flow_diagram}.png}
\end{figure}

\sphinxAtStartPar
In each subsequent section we break down the flows into individual components.


\subsubsection{1. Flow Out of Body}
\label{\detokenize{pages/tutorial_mc/addIO:flow-out-of-body}}
\sphinxAtStartPar
The flow out of the \sphinxstylestrong{Rest\_Of\_Body} is split into two flow rates as it moves
speices \sphinxstylestrong{A} and \sphinxstylestrong{B} to the \sphinxstylestrong{Liver\_Blood} and \sphinxstylestrong{Kidney}.  Below is the
part of the flow diagram we will be entering:

\begin{figure}[htbp]
\centering

\noindent\sphinxincludegraphics[scale=0.4]{{flow_diagram_rob}.png}
\end{figure}
\begin{enumerate}
\sphinxsetlistlabels{\arabic}{enumi}{enumii}{}{.}%
\item {} 
\sphinxAtStartPar
In \sphinxstylestrong{Options} select \sphinxstylestrong{Flow Between Compartments}. This is used to account
for flows leaving one compartment and entering one or more compartments.

\item {} 
\sphinxAtStartPar
Select the \sphinxstylestrong{Split Flow} checkbox. This option is selected when one outflow
is being split into multiple flows.  Keep \sphinxstylestrong{Number of Splits} to \sphinxstylestrong{2} as
we are spliting the flow to two compartments.

\item {} 
\sphinxAtStartPar
Set the information for the compartment the flow is leaving from. It has the
following specifications:
\begin{itemize}
\item {} 
\sphinxAtStartPar
\sphinxstylestrong{Flow\_Out\_of}: Rest\_Of\_Body

\item {} 
\sphinxAtStartPar
\sphinxstylestrong{Species Out}: A\_1

\item {} 
\sphinxAtStartPar
\sphinxstylestrong{Flow Variable}: F

\item {} 
\sphinxAtStartPar
\sphinxstylestrong{Flow Value}: 5

\end{itemize}

\sphinxAtStartPar
Current specifications need to have species in flow entered one at a time.
Meaning a new flow will be added for species \sphinxstylestrong{B}. With all paramters,
values can be changed at a later time in the parameter section.

\item {} 
\sphinxAtStartPar
Set the information from the first split flow into \sphinxstylestrong{Liver\_Blood}. Enter
the following:
\begin{itemize}
\item {} 
\sphinxAtStartPar
\sphinxstylestrong{Flow\_Out\_of}: Liver\_Blood

\item {} 
\sphinxAtStartPar
\sphinxstylestrong{Species Out}: A\_2

\item {} 
\sphinxAtStartPar
\sphinxstylestrong{Flow Variable}: F\_1

\item {} 
\sphinxAtStartPar
\sphinxstylestrong{Flow Value}: 3

\end{itemize}

\item {} 
\sphinxAtStartPar
Set the information from the second split flow into \sphinxstylestrong{Kidney}. Enter
the following:
\begin{itemize}
\item {} 
\sphinxAtStartPar
\sphinxstylestrong{Flow\_Out\_of}: Kidney

\item {} 
\sphinxAtStartPar
\sphinxstylestrong{Species Out}: A\_4

\item {} 
\sphinxAtStartPar
\sphinxstylestrong{Flow Variable}: F\_2

\item {} 
\sphinxAtStartPar
\sphinxstylestrong{Flow Value}: 2

\end{itemize}

\item {} 
\sphinxAtStartPar
Press the \sphinxstylestrong{Add} button to add this flow. You can select the
\sphinxstylestrong{Keep Active} checkbox on the bottom of the page to prevent this page
from closing on pressing the add button.

\begin{figure}[htbp]
\centering

\noindent\sphinxincludegraphics{{io_2_flow_between_A_marked}.png}
\end{figure}

\item {} 
\sphinxAtStartPar
Repeat the previous steps for species \sphinxstylestrong{B}. This should involve only
changing the \sphinxstylestrong{Species} dropdowns.
\begin{quote}

\begin{figure}[htbp]
\centering

\noindent\sphinxincludegraphics{{io_3_flow_between_B}.png}
\end{figure}
\end{quote}

\end{enumerate}

\begin{sphinxadmonition}{note}{Note:}
\sphinxAtStartPar
When adding \sphinxstylestrong{B}, we are reusing the flow variables \sphinxstylestrong{F}, \sphinxstylestrong{F\_1}, and
\sphinxstylestrong{F\_2}.  BioModME allows multiple saving of parameters but will check that
the parameters have the same base unit.  If not, the parameter will be
rejected. The value for the parameter will overwrite on each new overwrite
of the parameter.
\end{sphinxadmonition}


\subsubsection{2. Flow Back: Liver to Body}
\label{\detokenize{pages/tutorial_mc/addIO:flow-back-liver-to-body}}
\sphinxAtStartPar
The flow back from the liver to the rest of the body looks like:

\begin{figure}[htbp]
\centering

\noindent\sphinxincludegraphics[width=348\sphinxpxdimen,height=247\sphinxpxdimen]{{flow_diagram_l2rob}.png}
\end{figure}

\sphinxAtStartPar
This is like the previous step but with the compartments switched around.

\sphinxAtStartPar
Steps:
\begin{enumerate}
\sphinxsetlistlabels{\arabic}{enumi}{enumii}{}{.}%
\item {} 
\sphinxAtStartPar
Make sure \sphinxstylestrong{Options} is still on \sphinxstylestrong{Flow Between Compartments}.

\item {} 
\sphinxAtStartPar
Turn off \sphinxstylestrong{Split Flow}.

\item {} 
\sphinxAtStartPar
Set the information for the compartment the flow is leaving from. It has the
following specifications:
\begin{itemize}
\item {} 
\sphinxAtStartPar
\sphinxstylestrong{Flow Out Of}: Liver\_Blood

\item {} 
\sphinxAtStartPar
\sphinxstylestrong{Species Out}: A\_2

\item {} 
\sphinxAtStartPar
\sphinxstylestrong{Flow Variable}: F\_1

\item {} 
\sphinxAtStartPar
\sphinxstylestrong{Flow Value}: 3

\end{itemize}

\item {} 
\sphinxAtStartPar
Set the information for the compartment the flow is entering. It has the
following specifications:
\begin{itemize}
\item {} 
\sphinxAtStartPar
\sphinxstylestrong{Flow Into}: Rest\_Of\_Body

\item {} 
\sphinxAtStartPar
\sphinxstylestrong{Species In}: A\_1

\end{itemize}

\item {} 
\sphinxAtStartPar
Press the \sphinxstylestrong{Add} button to add this flow.
\begin{quote}

\begin{figure}[htbp]
\centering

\noindent\sphinxincludegraphics{{io_4_flow_back_A_lb_marked}.png}
\end{figure}
\end{quote}

\item {} 
\sphinxAtStartPar
Repeat the previous steps for species \sphinxstylestrong{B}. This should involve only
changing the \sphinxstylestrong{Species} dropdowns.
\begin{quote}

\begin{figure}[htbp]
\centering

\noindent\sphinxincludegraphics{{io_5_flow_back_b_lb}.png}
\end{figure}
\end{quote}

\end{enumerate}


\subsubsection{3. Flow Back: Kidney to Body}
\label{\detokenize{pages/tutorial_mc/addIO:flow-back-kidney-to-body}}
\sphinxAtStartPar
The flow back from the kidney to the rest of the body is:

\begin{figure}[htbp]
\centering

\noindent\sphinxincludegraphics[width=313\sphinxpxdimen,height=411\sphinxpxdimen]{{flow_diagram_k2rob}.png}
\end{figure}

\sphinxAtStartPar
Steps:
\begin{enumerate}
\sphinxsetlistlabels{\arabic}{enumi}{enumii}{}{.}%
\item {} 
\sphinxAtStartPar
Make sure \sphinxstylestrong{Options} is still on \sphinxstylestrong{Flow Between Compartments}.

\item {} 
\sphinxAtStartPar
Check that \sphinxstylestrong{Split Flow} is off.

\item {} 
\sphinxAtStartPar
Set the information for the compartment the flow is leaving from. It has the
following specifications:
\begin{itemize}
\item {} 
\sphinxAtStartPar
\sphinxstylestrong{Flow Out Of}: Kidney

\item {} 
\sphinxAtStartPar
\sphinxstylestrong{Species Out}: A\_4

\item {} 
\sphinxAtStartPar
\sphinxstylestrong{Flow Variable}: F\_2

\item {} 
\sphinxAtStartPar
\sphinxstylestrong{Flow Value}: 2

\end{itemize}

\item {} 
\sphinxAtStartPar
Set the information for the compartment the flow is entering. It has the
following specifications:
\begin{itemize}
\item {} 
\sphinxAtStartPar
\sphinxstylestrong{Flow Into}: Rest\_Of\_Body

\item {} 
\sphinxAtStartPar
\sphinxstylestrong{Species In}: A\_1

\end{itemize}

\item {} 
\sphinxAtStartPar
Press the \sphinxstylestrong{Add} button to add this flow.
\begin{quote}

\begin{figure}[htbp]
\centering

\noindent\sphinxincludegraphics{{io_6_flow_back_a_kidney_marked}.png}
\end{figure}
\end{quote}

\item {} 
\sphinxAtStartPar
Repeat the previous steps for species \sphinxstylestrong{B}. This should involve only
changing the \sphinxstylestrong{Species} dropdowns.
\begin{quote}

\begin{figure}[htbp]
\centering

\noindent\sphinxincludegraphics{{io_7_flow_back_b_kidney}.png}
\end{figure}
\end{quote}

\end{enumerate}


\subsubsection{4. Clearance of A From Kidney}
\label{\detokenize{pages/tutorial_mc/addIO:clearance-of-a-from-kidney}}
\sphinxAtStartPar
Drug A is excreted from the kidney at a constant rate. The isolated process is
shown below:

\begin{figure}[htbp]
\centering

\noindent\sphinxincludegraphics[width=405\sphinxpxdimen,height=275\sphinxpxdimen]{{flow_clearance}.png}
\end{figure}

\sphinxAtStartPar
Steps:
\begin{enumerate}
\sphinxsetlistlabels{\arabic}{enumi}{enumii}{}{.}%
\item {} 
\sphinxAtStartPar
Select \sphinxstylestrong{Clearance} in the \sphinxstylestrong{Options} dropdown.

\item {} 
\sphinxAtStartPar
Enter the following information in the main box:
\begin{itemize}
\item {} 
\sphinxAtStartPar
\sphinxstylestrong{Compartment}: Kidney

\item {} 
\sphinxAtStartPar
\sphinxstylestrong{Species}: B\_4

\item {} 
\sphinxAtStartPar
\sphinxstylestrong{Rate}: k\_e

\end{itemize}

\item {} 
\sphinxAtStartPar
Press the \sphinxstylestrong{Add} button to add the clearance of B from the kidney.

\end{enumerate}

\begin{figure}[htbp]
\centering

\noindent\sphinxincludegraphics{{io_8_clearance_marked}.png}
\end{figure}


\subsubsection{5. Facilitated Diffusion of A}
\label{\detokenize{pages/tutorial_mc/addIO:facilitated-diffusion-of-a}}
\sphinxAtStartPar
The next two seconds will look at diffusion processes from the liver blood
to the liver tissue and back.  Below is the facilitated diffusion of molecule
\sphinxstylestrong{A} from the liver blood to the liver tissue.

\begin{figure}[htbp]
\centering

\noindent\sphinxincludegraphics[width=487\sphinxpxdimen,height=241\sphinxpxdimen]{{flow_facilitated_diffusion}.png}
\end{figure}

\sphinxAtStartPar
Steps:
\begin{enumerate}
\sphinxsetlistlabels{\arabic}{enumi}{enumii}{}{.}%
\item {} 
\sphinxAtStartPar
Select \sphinxstylestrong{Facilitated Diffusion} in the \sphinxstylestrong{Options} dropdown.

\item {} 
\sphinxAtStartPar
Enter the following information in the first row of the main box:
\begin{itemize}
\item {} 
\sphinxAtStartPar
\sphinxstylestrong{From Compartment}: Liver\_Blood

\item {} 
\sphinxAtStartPar
\sphinxstylestrong{From Species}: A\_2

\item {} 
\sphinxAtStartPar
\sphinxstylestrong{Vmax}: fVmax\_1

\item {} 
\sphinxAtStartPar
\sphinxstylestrong{Km}: fKm\_1

\end{itemize}

\item {} 
\sphinxAtStartPar
Enter the following information in the second row of the main box:
\begin{itemize}
\item {} 
\sphinxAtStartPar
\sphinxstylestrong{To Compartment}: Liver\_Tissue

\item {} 
\sphinxAtStartPar
\sphinxstylestrong{From Species}: A\_3

\end{itemize}

\item {} 
\sphinxAtStartPar
Press the \sphinxstylestrong{Add} button to add this filitated diffusion flow to the model.

\end{enumerate}

\begin{figure}[htbp]
\centering

\noindent\sphinxincludegraphics{{io_9_facdif_marked}.png}
\end{figure}


\subsubsection{6. Simple Diffusion of B}
\label{\detokenize{pages/tutorial_mc/addIO:simple-diffusion-of-b}}
\begin{figure}[htbp]
\centering

\noindent\sphinxincludegraphics[width=487\sphinxpxdimen,height=241\sphinxpxdimen]{{flow_simple_diffusion}.png}
\end{figure}

\sphinxAtStartPar
Steps:
\begin{enumerate}
\sphinxsetlistlabels{\arabic}{enumi}{enumii}{}{.}%
\item {} 
\sphinxAtStartPar
Select \sphinxstylestrong{Simple Diffusion} in the \sphinxstylestrong{Options} dropdown.

\item {} 
\sphinxAtStartPar
Enter the following information in the first row of the main box:
\begin{itemize}
\item {} 
\sphinxAtStartPar
\sphinxstylestrong{Compartment}: Liver\_Blood

\item {} 
\sphinxAtStartPar
\sphinxstylestrong{Species}: B\_2

\item {} 
\sphinxAtStartPar
\sphinxstylestrong{Diffusivity Coefficient}: PS

\end{itemize}

\item {} 
\sphinxAtStartPar
Enter the following information in the second row of the main box:
\begin{itemize}
\item {} 
\sphinxAtStartPar
\sphinxstylestrong{Compartment}: Liver\_Tissue

\item {} 
\sphinxAtStartPar
\sphinxstylestrong{Species}: B\_3

\end{itemize}

\begin{sphinxadmonition}{note}{Note:}
\sphinxAtStartPar
The order of entered compartments and species does not matter for
simple diffusion.
\end{sphinxadmonition}

\item {} 
\sphinxAtStartPar
Press the \sphinxstylestrong{Add} button to add this simple diffusion flow to the model.

\end{enumerate}

\begin{figure}[htbp]
\centering

\noindent\sphinxincludegraphics{{io_10_simpdiff_marked}.png}
\end{figure}

\sphinxAtStartPar
This should we the last term entered in the Input/Output box. There should be
nine terms in the results table.


\subsection{Parameters}
\label{\detokenize{pages/tutorial_mc/parameters:parameters}}\label{\detokenize{pages/tutorial_mc/parameters::doc}}
\sphinxAtStartPar
Navigate to the next box: \sphinxstylestrong{Parameters}.

\sphinxAtStartPar
Below is the generated parameter table that has been created so far by the
above entered information. All descriptions are pregenerated based on where
the parameter originated. Like other tables, all values are editable.

\begin{figure}[htbp]
\centering

\noindent\sphinxincludegraphics{{param_0}.png}
\end{figure}

\sphinxAtStartPar
The parameter values for the model are as follows:

\begin{figure}[htbp]
\centering

\noindent\sphinxincludegraphics[scale=0.5]{{parameter_table}.png}
\end{figure}

\sphinxAtStartPar
The below figure shows the parameter table with the values entered.

\begin{figure}[htbp]
\centering

\noindent\sphinxincludegraphics{{param_1}.png}
\end{figure}

\begin{sphinxadmonition}{note}{Note:}
\sphinxAtStartPar
In the Figure above, we did not keep the units of our parameters
consistent. In some cases, ml was used and in others liters was used.
This is okay.  The application converts all units to a “base” unit
in the background and performs all calculations at that base level.
\end{sphinxadmonition}


\subsection{Differential Equations}
\label{\detokenize{pages/tutorial_mc/differentialEquations:differential-equations}}\label{\detokenize{pages/tutorial_mc/differentialEquations::doc}}
\sphinxAtStartPar
Navigate to the final box \sphinxstylestrong{“Differential Equations”}.

\sphinxAtStartPar
This tab is used to generate and see the mathematical equations that govern
our model.  This tab will begin with all species being equal to “NA” until the
generate button is pressed or model is executed.

\sphinxAtStartPar
Press the “Generate” button to solve the system for its differential equation
output and display them in text form.

\begin{figure}[htbp]
\centering

\noindent\sphinxincludegraphics{{diffeqn_1}.png}
\end{figure}

\sphinxAtStartPar
\sphinxstylestrong{Note:} The right sidebar of the “System of Differential Equations” box has an
option for viewing the differential equations including turning off the math
rendering and showing each derivation line as a newline for each differential
equation.


\section{Results}
\label{\detokenize{pages/tutorial_mc/2_model_results:results}}\label{\detokenize{pages/tutorial_mc/2_model_results::doc}}
\sphinxAtStartPar
In this section, we will cover exectuting the model and its different outputs.


\subsection{Solve Model}
\label{\detokenize{pages/tutorial_mc/2_model_results:solve-model}}
\sphinxAtStartPar
Navigate to the “Execute Model” tab on the lefthand tabbar. This tab solves our
model and provides us with the options to customize the solver. For this model,
we will just solve using the standard options:
\begin{enumerate}
\sphinxsetlistlabels{\arabic}{enumi}{enumii}{}{.}%
\item {} 
\sphinxAtStartPar
Enter the following times: Starting = 0, End = 5, Step = 0.001, Unit = hr.
This generates the model to be solved and integrated at time points 0, 0.001,
0.002…..5.998, 5.999, 6.

\item {} 
\sphinxAtStartPar
Press the “Run Solver” button.

\item {} 
\sphinxAtStartPar
An output table should be generated with the concentration of each species
along each time step of the solved model.

\end{enumerate}

\sphinxAtStartPar
You can download this solved model data to a csv file in the
download tab on this page.  Post processing is not used in this tutorial.

\begin{figure}[htbp]
\centering

\noindent\sphinxincludegraphics{{execute_1}.png}
\end{figure}

\begin{sphinxadmonition}{note}{Note:}
\sphinxAtStartPar
BioModMe is set by default to show the results table rounded to 3
digits. Most terms in these results at the beginning will round
to 0. To change this, check the \sphinxstylestrong{Viewing Options} box and turn
off the \sphinxstylestrong{Round} option or turn on the \sphinxstylestrong{Scientific Notation}
option.
\end{sphinxadmonition}

\sphinxAtStartPar
Below is how the table for this model should look after the initial solve:

\begin{figure}[htbp]
\centering

\noindent\sphinxincludegraphics{{results_table}.png}
\end{figure}

\sphinxAtStartPar
To see the values for the model in a good form, we will change the viewing
options:
\begin{enumerate}
\sphinxsetlistlabels{\arabic}{enumi}{enumii}{}{.}%
\item {} 
\sphinxAtStartPar
Click the \sphinxstylestrong{Viewing Options} Checkbox.

\item {} 
\sphinxAtStartPar
Change Results Units to \sphinxstylestrong{umol}.

\end{enumerate}

\begin{figure}[htbp]
\centering

\noindent\sphinxincludegraphics{{result_table_viewing_options}.png}
\end{figure}

\sphinxAtStartPar
In this section, we will examine the plotting of this model.


\subsection{Visualization}
\label{\detokenize{pages/tutorial_mc/2_model_results:visualization}}
\sphinxAtStartPar
Move to the \sphinxstylestrong{Visualization} tab.  The starting plot from the model should
look like the following:

\begin{figure}[htbp]
\centering

\noindent\sphinxincludegraphics{{plot_1}.png}
\end{figure}

\sphinxAtStartPar
For this problem we want to examine how the concentration of \sphinxstylestrong{A} varies
throughout the bloodsteam (\sphinxstylestrong{A\_1}, \sphinxstylestrong{A\_2}, \sphinxstylestrong{A\_4}). Click the \sphinxstylestrong{Variable}
dropdown and remove all necessary variables:

\begin{figure}[htbp]
\centering

\noindent\sphinxincludegraphics{{plot_var}.png}
\end{figure}

\sphinxAtStartPar
The resulting plot should look like:

\begin{figure}[htbp]
\centering

\noindent\sphinxincludegraphics{{plot_2}.png}
\end{figure}


\chapter{Parameter Estimation Tutorial}
\label{\detokenize{pages/tutorial_par_estimation:parameter-estimation-tutorial}}\label{\detokenize{pages/tutorial_par_estimation::doc}}
\sphinxAtStartPar
This is a basic parameter estimation tutorial. To download the example data
used in this tutorial go to
\sphinxurl{https://github.com/Jwomack7512-bio/BobTheBuilder/blob/main/shared\_files/Nicotine\_cigarette.csv}.

\sphinxAtStartPar
We assume you have basic knowledge of creating a model and will skip showing
every step of that process. Please see the
{\hyperref[\detokenize{pages/backend/parameter_estimation::doc}]{\sphinxcrossref{\DUrole{doc}{parameter estimation documentation}}}}
for details on the solver and options available to use in the solver.

\sphinxAtStartPar
We will construct the following simple three compartment model:

\begin{figure}[htbp]
\centering

\noindent\sphinxincludegraphics{{0_diagram_and_table}.png}
\end{figure}

\sphinxAtStartPar
In the above diagram, we model the concentration of a drug, Nicotine (\sphinxstylestrong{C}), as
it moves from the alveoli of the lungs, through the plasma of the body, and a
compartment summarizing various tissues of the body. We assume Nicotine flows
out of the alveoli with flow rate (\sphinxstylestrong{F\_1}), into the plasma of the body.
From there, Nicotine diffuses to various tissues of the body through simple
diffusion (\sphinxstylestrong{PS}) and is also removed from the plasma at a flow rate of
\sphinxstylestrong{F\_2}.  In this model, we are unsure of the value of \sphinxstylestrong{PS} and will use
parameter estimation to estimate the value from experimental data that found
the concentration of nicotine in the plasma at various timepoints. The starting
concentration of \sphinxstylestrong{C} in the alveoli is 1.2 mol.


\section{Build Model}
\label{\detokenize{pages/tutorial_par_estimation:build-model}}
\sphinxAtStartPar
To being buiding the model, we will add our three compartments with their
corresponding volumes:
\begin{itemize}
\item {} 
\sphinxAtStartPar
Compartment: Alveoli, Volume = 1 L

\item {} 
\sphinxAtStartPar
Compartment: Plasma, Volume = 7.8 L

\item {} 
\sphinxAtStartPar
Compartment: Various\_Tissues, Volume = 9.5 L

\end{itemize}

\begin{figure}[htbp]
\centering

\noindent\sphinxincludegraphics{{1_comps_named}.png}
\end{figure}

\sphinxAtStartPar
Next, we will add a concentration (\sphinxstylestrong{C}) variable for nicotine to each
compartment and give an initial concentration of 1.2 mol to \sphinxstylestrong{C\_1}.

\begin{figure}[htbp]
\centering

\noindent\sphinxincludegraphics{{2_species_add_w_con}.png}
\end{figure}

\sphinxAtStartPar
Next, we need to add the compartment input/outputs.  We have three to add:
\begin{enumerate}
\sphinxsetlistlabels{\arabic}{enumi}{enumii}{}{.}%
\item {} 
\sphinxAtStartPar
Flow between Alveoli and Plasma

\begin{figure}[htbp]
\centering

\noindent\sphinxincludegraphics{{3_flow_1_add}.png}
\end{figure}

\item {} 
\sphinxAtStartPar
Flow out of Plasma

\begin{figure}[htbp]
\centering

\noindent\sphinxincludegraphics{{3_flow_2_add}.png}
\end{figure}

\item {} 
\sphinxAtStartPar
Simple diffusion between Plasma and Various\_Tissues

\begin{figure}[htbp]
\centering

\noindent\sphinxincludegraphics{{3_flow_PS_add}.png}
\end{figure}

\end{enumerate}

\sphinxAtStartPar
With all our model components constructed, we need to assign parameter values.
We do not know the value of PS but we will give the model a reasonable
starting estimate (0.5 L/min).

\begin{figure}[htbp]
\centering

\noindent\sphinxincludegraphics{{4_parameters}.png}
\end{figure}

\begin{sphinxadmonition}{note}{Note:}
\sphinxAtStartPar
The parameter table rounds to two decimal places.  This is why F\_2
shows as 0.08 in the table. If you double click the value it will
show as the appropriate 0.078.
\end{sphinxadmonition}

\sphinxAtStartPar
For reference, here are the solved differential equations:

\begin{figure}[htbp]
\centering

\noindent\sphinxincludegraphics{{5_diff_eqn}.png}
\end{figure}

\sphinxAtStartPar
Move to the \sphinxstylestrong{Execute Model} tab and enter the time for this model:
\begin{itemize}
\item {} 
\sphinxAtStartPar
Start Time = 0

\item {} 
\sphinxAtStartPar
End Time = 70

\item {} 
\sphinxAtStartPar
Time Step = 1

\item {} 
\sphinxAtStartPar
Unit = min

\end{itemize}

\begin{figure}[htbp]
\centering

\noindent\sphinxincludegraphics{{6_execute}.png}
\end{figure}

\sphinxAtStartPar
The beginning of the results table should look like this:

\begin{figure}[htbp]
\centering

\noindent\sphinxincludegraphics{{6_execute_table}.png}
\end{figure}


\section{Visualize Model}
\label{\detokenize{pages/tutorial_par_estimation:visualize-model}}
\sphinxAtStartPar
Move to the \sphinxstylestrong{Visualization} tab. A starting plot should be generated that
looks like:

\begin{figure}[htbp]
\centering

\noindent\sphinxincludegraphics{{7_results_plot}.png}
\end{figure}

\sphinxAtStartPar
We want to examine the the concentration of Nicotine in the plasma (\sphinxstylestrong{C\_2}).
Go into the \sphinxstylestrong{Variables} dropdown and remove variables \sphinxstylestrong{C\_1} and \sphinxstylestrong{C\_3}.

\begin{figure}[htbp]
\centering

\noindent\sphinxincludegraphics{{7_results_plot_c2}.png}
\end{figure}

\sphinxAtStartPar
We should import our data for Nicotine in the plasma to overlay with the plot.
Scroll below the plot to the box header that is titled \sphinxstylestrong{Import Data}.
Browse your desktop for the nictoine dataset.  Which can be downloaded from
our Github. Be sure to click the \sphinxstylestrong{Apply Overlay} checkbox to show the imported
data. For importing data, the first row of the data is the headers. The first
column should be time while the remaining columns should be your concentrations
at those times. The header names have to match the model names for the data to
be correctly plotted.

\begin{figure}[htbp]
\centering

\noindent\sphinxincludegraphics{{8_import_data}.png}
\end{figure}

\sphinxAtStartPar
The resulting plot should look like the figure below:

\begin{figure}[htbp]
\centering

\noindent\sphinxincludegraphics{{8_import_data_plot}.png}
\end{figure}

\sphinxAtStartPar
In this overlayed plot, we can see our model has the general shape of our data
but doesn’t quite mimic it well. We will run parameter estimation for \sphinxstylestrong{PS} to
see if we can find a better fit.


\section{Perform Parameter Estimation}
\label{\detokenize{pages/tutorial_par_estimation:perform-parameter-estimation}}
\sphinxAtStartPar
Go to the \sphinxstylestrong{Modeler’s Toolbox} Tab, subtab \sphinxstylestrong{Parameter Estimation}. Here, we
will start by importing our data into this module.

\begin{figure}[htbp]
\centering

\noindent\sphinxincludegraphics{{9_pe_import_data}.png}
\end{figure}

\sphinxAtStartPar
Move to the next box on the page \sphinxstylestrong{Select Parameters}. Use the dropdown to
select \sphinxstylestrong{PS}.

\begin{figure}[htbp]
\centering

\noindent\sphinxincludegraphics{{9_pe_table}.png}
\end{figure}

\sphinxAtStartPar
There are five columns in the generated table. They are as follows:
\begin{itemize}
\item {} 
\sphinxAtStartPar
Parameters \sphinxhyphen{} the selected values to estimate.

\item {} 
\sphinxAtStartPar
Initial Guess \sphinxhyphen{} the starting value to begin parameter estimation at.

\item {} 
\sphinxAtStartPar
Lower Bound \sphinxhyphen{} the lowest acceptable value this parameter can be.

\item {} 
\sphinxAtStartPar
Upper Bound \sphinxhyphen{} the highest acceptable value this parameter can be.

\item {} 
\sphinxAtStartPar
Calculated Value \sphinxhyphen{} the found value from parameter estimation after calculations.

\end{itemize}

\sphinxAtStartPar
Values for Lower and Upper bound can be left blank if no bounds want to be used.
Here we use the following:
\begin{itemize}
\item {} 
\sphinxAtStartPar
Initial Guess \sphinxhyphen{} 0.50

\item {} 
\sphinxAtStartPar
Lower Bound \sphinxhyphen{} 0

\item {} 
\sphinxAtStartPar
Upper Bound \sphinxhyphen{} 1

\end{itemize}

\sphinxAtStartPar
Press the \sphinxstylestrong{Run} button and the program should output a value of approximately
0.10.

\begin{figure}[htbp]
\centering

\noindent\sphinxincludegraphics{{9_pe_table_results}.png}
\end{figure}

\sphinxAtStartPar
The next box \sphinxstylestrong{Estimation Results} will contain the model fit of the variable
with its corrsponding data along with the results of the iterations of the
parameter estimation algorithm.

\begin{figure}[htbp]
\centering

\noindent\sphinxincludegraphics{{10_pe_plot_results}.png}
\end{figure}

\sphinxAtStartPar
Press the \sphinxstylestrong{Store} button to overwrite your current parameters in the model
with the estimated values.

\begin{sphinxadmonition}{note}{Note:}
\sphinxAtStartPar
You can estimate as many parameters as you want with this setup.
Just note that the more uncertainy you introduce into your model the
longer the algorithm can take to find a solution. It can also affect
the probability of finding a solution.
\end{sphinxadmonition}


\chapter{Welcome to BioModME’s documentation!}
\label{\detokenize{index:welcome-to-biomodme-s-documentation}}
\sphinxAtStartPar
BioSysMod is an application that is meant to streamline building biological
computational models. It is an all\sphinxhyphen{}in\sphinxhyphen{}one tool that allows the user to build
a model, solve for its mathematical equations, and plot all relevant features.


\chapter{First Steps}
\label{\detokenize{index:first-steps}}
\sphinxAtStartPar
Are you new to BioModME or computational modeling? This is the place to start!
\begin{itemize}
\item {} 
\sphinxAtStartPar
\sphinxstylestrong{Access:}
{\hyperref[\detokenize{pages/gettingStarted/useonline::doc}]{\sphinxcrossref{\DUrole{doc}{Online Application}}}} |
{\hyperref[\detokenize{pages/gettingStarted/downloadtoR::doc}]{\sphinxcrossref{\DUrole{doc}{Download}}}}

\item {} 
\sphinxAtStartPar
\sphinxstylestrong{Tutorials:}
{\hyperref[\detokenize{pages/tutorial_sc_index::doc}]{\sphinxcrossref{\DUrole{doc}{Basic Model}}}} |
{\hyperref[\detokenize{pages/tutorial_mc_index::doc}]{\sphinxcrossref{\DUrole{doc}{MultiCompartment Model}}}}

\end{itemize}


\chapter{How the documentation is organized}
\label{\detokenize{index:how-the-documentation-is-organized}}
\sphinxAtStartPar
The documentation is split into the following sections:
\begin{description}
\item[{\sphinxstylestrong{Getting Started}}] \leavevmode
\sphinxAtStartPar
Contains all the information on accessing/downloading the application for
use.

\item[{\sphinxstylestrong{API}}] \leavevmode
\sphinxAtStartPar
Explains the specifics of tabs and pages of the application. If you are
wondering what a specific option or button does, this is the place to check!

\item[{\sphinxstylestrong{Mathematical Information}}] \leavevmode
\sphinxAtStartPar
This section contains all the mathematical explainations that goes into
solving the mathematical differential equations. See how the equations are
derived or what a specific equation building option results in.

\item[{\sphinxstylestrong{Tutorials}}] \leavevmode
\sphinxAtStartPar
Great place to start to learn the ins and outs of how to use BioModME. We
provide tutorials for multiple model types and tutorials on how to use some
of the more tricky features of the application.

\end{description}


\chapter{Links}
\label{\detokenize{index:links}}\begin{itemize}
\item {} 
\sphinxAtStartPar
Source Code (stable): \sphinxurl{https://github.com/MCWComputationalBiologyLab/BioModME}

\item {} 
\sphinxAtStartPar
Source Code (dev): \sphinxurl{https://github.com/Jwomack7512-bio/BobTheBuilder}

\item {} 
\sphinxAtStartPar
Home Page: \sphinxurl{https://biomodme.ctsi.mcw.edu/}

\item {} 
\sphinxAtStartPar
Contact: \sphinxhref{mailto:jwomack@mcw.edu}{jwomack@mcw.edu}, \sphinxhref{mailto:rdash@mcw.edu}{rdash@mcw.edu}

\end{itemize}



\renewcommand{\indexname}{Index}
\printindex
\end{document}